% Routing outcome model --- flow-aware, officer-observable
% Branch: muse/routing-experiments. Experimental; not near main.
% Audience: first-year graduate econometrics.
% Body = the routing problem. Appendices = foundations, proved.
% Prerequisites are taken from ../probability (Mohanty, "Basic mathematics for
% statistics and economics"); Appendix A says exactly what is assumed.
% Compile: latexmk -pdf routing-outcome-model.tex
\documentclass[11pt,a4paper]{article}
\usepackage[margin=1in]{geometry}
\usepackage{booktabs,amsmath,amssymb,amsthm,hyperref,graphicx,xcolor,longtable,enumitem}
\hypersetup{colorlinks=true,linkcolor=blue,citecolor=blue,urlcolor=blue}

\theoremstyle{plain}
\newtheorem{theorem}{Theorem}[section]
\newtheorem{lemma}[theorem]{Lemma}
\newtheorem{proposition}[theorem]{Proposition}
\newtheorem{corollary}[theorem]{Corollary}
\theoremstyle{definition}
\newtheorem{assumption}{Assumption}
\newtheorem{definition}[theorem]{Definition}
\newtheorem{example}[theorem]{Example}
\theoremstyle{remark}
\newtheorem*{remark*}{Remark}

\DeclareMathOperator*{\argmin}{arg\,min}
\DeclareMathOperator*{\argmax}{arg\,max}
\DeclareMathOperator{\Var}{Var}
\DeclareMathOperator{\Cov}{Cov}
\DeclareMathOperator{\Corr}{Corr}
\newcommand{\E}{\mathbb E}
\newcommand{\Prob}{\mathbb P}
\newcommand{\indep}{\perp\!\!\!\perp}
\newcommand{\one}[1]{\mathbf 1\{#1\}}
\newcommand{\code}[1]{\texttt{\small #1}}
\newcommand{\notbuilt}{\textcolor{red}{\textbf{[not built]}}}
\newcommand{\probrepo}{\textsc{bmse}}

\title{Outcome-Based Routing of Public Grievances\\
\large Minimising disposal time subject to a correctness constraint:\\
estimand, identification, and estimation}
\author{Janasunani 2.0 --- DPIC}
\date{13 August 2026}

\begin{document}
\maketitle

\begin{abstract}
A citizen files a grievance; an officer decides where to send it; some weeks or months later
it closes. This document asks whether that routing decision can be made better --- faster
disposal, no loss of correctness --- using only information the officer already has on screen
at the moment of assignment.

The body is about the routing problem. Section~\ref{sec:setup} fixes notation and gets the
outcome definition right: closing a grievance \emph{correctly} and \emph{helping the
petitioner} are different events, and both differ again from the case merely being closed.
Section~\ref{sec:estimand} converts the policy question into a constrained optimisation over
decision rules, and explains why the natural first formulation --- ``time to disposal among
correctly disposed cases'' --- conditions on a post-treatment variable and must be replaced.
Section~\ref{sec:ident} states the identifying assumptions, argues for each in this
institutional setting, and assembles the estimating equation. Section~\ref{sec:est} builds the
estimator in five layers and says which statistical model estimates which nuisance function,
and why. Sections~\ref{sec:policy}--\ref{sec:interfere} cover policy learning, diagnostics,
sensitivity, and the queueing interference a routing policy necessarily creates.
Section~\ref{sec:evidence} reports what has been computed.

The general theory is in the appendices, stated and proved.
Appendix~\ref{app:prereq} lists exactly what is taken as known from the companion \probrepo{}
reference, and what is not. Appendix~\ref{app:semipar} sets out the semiparametric vocabulary
--- models, parameters as functionals, target versus nuisance, influence functions, the
efficiency bound, Neyman orthogonality --- on which everything after it rests.
Appendix~\ref{app:causal} proves causal identification and double robustness,
Appendix~\ref{app:censor} handles censoring, Appendix~\ref{app:estlemmas} collects the
estimation lemmas, and Appendix~\ref{app:sens} develops the sensitivity model.
\end{abstract}

\tableofcontents
\newpage

%=====================================================================
\section{The problem}\label{sec:intro}
%=====================================================================

\subsection{Institutional setting}

Odisha's Janasunani system receives citizen grievances --- a pension not disbursed, a ration
card refused, a road not repaired, a boundary dispute. Each is created in an electronic
register, assigned to an office, possibly forwarded along a chain of officers, and eventually
closed with a disposal remark. The register records the creation date, the closure date, the
text of the closing remark, a flag for whether the beneficiary benefited, and the full chain
of officers who held the file.

Routing today is historical. The production rule sends a grievance to
\[
\argmax_{d}\ \Prob(\text{department}=d \mid \text{category}, \text{district}),
\]
that is, to wherever cases of this type have gone before. This has a specific blind spot: it
learns where cases \emph{were} sent, never where they were \emph{handled well}. A department
that receives every land dispute in a district and resolves none of them is, to this rule,
exactly the right destination.

\subsection{Why the obvious objective is wrong}

The first instinct is to minimise time to closure. This fails, instructively.

Closure is an action the officer controls directly. A grievance can be closed at any moment by
recording that it has been disposed. If the objective rewards short durations without
qualification, the optimal behaviour under that objective is to close every case immediately
without doing anything. This is not hypothetical: the most common closing remark in the system
records disposal with no action claimed.

The objective must therefore be conditional in some way on the work having been done. Making
that precise is most of the content of \S\ref{sec:outcome} and \S\ref{sec:estimand}. The short
version: we minimise expected duration \emph{subject to a constraint} that correctness not
deteriorate, rather than minimising duration \emph{among} the cases that came out correct.
That is the difference between a well-posed causal question and one that conditions on the
future.

\subsection{Why this setting is unusually favourable}

Routing studies are usually hopeless observationally, because the assigner knows things the
analyst does not, and those things drive both assignment and outcome. Here the situation is
better for a concrete institutional reason: the assigning officer works from a screen, and we
have the screen.

The officer sees the grievance text, the category and subcategory, the district and block, the
mode of filing, whether documents are attached, the scheme tag, the current backlog of the
candidate destination, and that destination's recent performance. All of it is in the
database. This does not make unconfoundedness true, but it changes what one must believe for
it to be true: not ``no unobserved case characteristics'' but ``no unwritten judgement''.
Section~\ref{sec:ident} makes this precise and \S\ref{sec:sens} quantifies how much unwritten
judgement would be needed to overturn a conclusion.

%=====================================================================
\section{Setup}\label{sec:setup}
%=====================================================================

\subsection{Units and sampling}

The unit is a grievance, indexed $i=1,\dots,n$, arriving at calendar time $t_{0i}$. We treat
grievances as a sample from a population distribution $P$. The sample is not i.i.d.\ in a
strict sense --- grievances arriving in the same district in the same period share shocks ---
and \S\ref{sec:variance} handles that through clustering rather than by pretending otherwise.

\subsection{Covariates}\label{sec:covariates}

$X_i \in \mathcal X$ collects everything observable at $t_{0i}$, of four kinds:

\begin{enumerate}[leftmargin=*]
\item \textbf{Case characteristics}: category, subcategory, district, block, mode of filing,
whether supporting documents are attached, the redacted grievance text, a disability flag, the
scheme tag.
\item \textbf{The citizen's intake choice}: the coarse office selected when filing. This is
\emph{not} the treatment --- the citizen chooses it, not the officer --- but it predicts both
routing and outcome, so it belongs in $X$.
\item \textbf{Congestion}: the backlog $Q_r(t_{0i})$ of each candidate destination role $r$ at
the moment of assignment. This is the number on the officer's dashboard, and
\S\ref{sec:backlog} constructs it.
\item \textbf{Trailing performance}: for each candidate destination, its median duration and
action rate over the preceding 90 days, computed strictly from cases closed before $t_{0i}$.
\end{enumerate}

The point-in-time restriction on items 3 and 4 is not a technicality. A feature computed from
the future would let the model use information the officer could not have had, which breaks
both the interpretation --- we are evaluating a rule the officer could have followed --- and
the identification, since \S\ref{sec:ident} requires $X$ to be pre-treatment.

\subsubsection{Measuring the backlog}\label{sec:backlog}

``Backlog'' is ambiguous and the obvious implementations are wrong in instructive ways.

\paragraph{What we want.} $Q_r(t)$ should be the number the officer would see on looking at
role $r$'s queue at instant $t$: files that have arrived with $r$ and not yet left. Two
quantifiers need pinning down. \emph{Whose} queue --- an individual officer's or the role's?
We use the role, because the treatment is a role sequence (Definition~\ref{def:flow}) and an
individual officer's identity is not something the policy selects. And \emph{left} how ---
closed, or forwarded onward? Either departure frees the officer, so both count.

\paragraph{Custody intervals.} The register stores no queue length, so it must be
reconstructed. The raw material is the action history, which records for each event the
ticket, the timestamp, and the acting user. Mapping the acting user to a role gives, for each
ticket, a sequence of custody events; between consecutive events the ticket sits with one
role. For ticket $j$ with custody events at times $u_{j,1}<u_{j,2}<\dots$ held by roles
$r_{j,1},r_{j,2},\dots$, define the half-open custody interval
\[
I_{j,k} = \bigl[\,u_{j,k},\ \min(u_{j,k+1},\ \text{closure}_j)\,\bigr),
\]
with the first interval beginning at the assignment date rather than the first action. Then
\begin{equation}
Q_r(t) = \#\bigl\{(j,k) : r_{j,k}=r,\ t\in I_{j,k}\bigr\}.
\label{eq:Q}
\end{equation}
Capacity is trailing throughput, $\text{cap}_r(t)=\#\{$departures from role $r$ in
$[t-90,t)\}$, so $Q_r(t)/\text{cap}_r(t)$ has units of days of work outstanding --- the
quantity entering \eqref{eq:micro}, and comparable across roles of very different size.

\paragraph{Two traps.} First, \eqref{eq:Q} must be evaluated using only events dated before
$t$. The membership test $t\in I_{j,k}$ looks like it consults the future, since the interval's
right endpoint may lie after $t$; it does not, because the officer at $t$ needs to know only
that \emph{no departure has happened yet}, which is observable at $t$. The equivalent
backward-looking form --- arrived before $t$, no departure event before $t$ --- is the one to
implement and the one to state when defending the point-in-time claim.

Second, and more dangerous: a naive implementation writes the condition as
\code{assigned\_on <= t AND t < resolved\_on}. For a grievance never resolved,
\code{resolved\_on} is null, the comparison evaluates to null rather than true, and the ticket
silently drops out. Since unresolved tickets \emph{are} the backlog, this undercounts $Q$
exactly where $Q$ is largest, and worst in the periods of heaviest congestion. The condition
must treat a null closure as ``still open''.

\paragraph{Status.} $Q_r(t)$ and the trailing-performance features are specified here and
\notbuilt. The implemented design matrix has case characteristics, the intake office, the
flow-derived columns and time effects, but no congestion measure and no trailing performance.
This is not cosmetic: \S\ref{sec:ident} justifies unconfoundedness on the ground that $X$
reproduces the officer's screen, and congestion is on the screen. It is also a textbook
confounder --- a busy office is both slower and plausibly routed to differently --- so its
omission biases the flow comparison in an unknown direction.

\subsection{Treatment: the flow}\label{sec:flow}

The register contains a field recording the citizen's intake office --- eight coarse values
such as ``Collector'' or ``Office of Chief Minister''. It separately contains
\code{vchAllEscUser}, a comma-separated list of the internal identifiers of every officer who
held the file. These are not the same thing and not even close: of the 693,691 grievances filed
at ``Collector'', only 12.8\% are ever actually pending with a Collector. A policy that
reranked the intake label would be optimising the citizen's guess.

\begin{definition}[Flow]\label{def:flow}
Decode each user identifier in \code{vchAllEscUser} to that user's role, via the user--role and
role--name master tables. The \emph{flow} of grievance $i$ is the resulting role sequence,
written as a template string: \code{6,8} denotes Secretary$\to$Chief Minister's Office,
\code{1,5} denotes BDO$\to$Collector.
\end{definition}

The treatment is $A_i\in\mathcal A$, the flow.

\emph{Why roles and not users.} There are 2,747 distinct user identifiers in the chains.
Treating each as a treatment level would make every comparison a comparison of individuals,
with no overlap and no hope of positivity. Collapsing to roles gives 1,318 distinct sequences,
a few hundred with meaningful support, while preserving what differs between routings: which
\emph{kind} of officer handles the case, and how many hands it passes through.

\emph{Why the whole sequence.} Each additional link is a physical handoff: the file is
re-queued, re-read, re-prioritised. Sequence length is a treatment component with a direct
mechanical effect on duration, not merely a label; \S\ref{sec:micro} models it.

\begin{definition}[Admissible set]\label{def:admissible}
$\mathcal A(x)\subseteq\mathcal A$ is the set of flows assignable to a grievance with
covariates $x$: the intersection of (i) \emph{support} --- the flow was used at least
$\underline m$ times for this category--district cell in the training period; (ii)
\emph{structural feasibility} --- each consecutive role pair is an edge in the forwarding graph
derived from the escalation master tables, so the routing is administratively permitted; (iii)
\emph{jurisdiction} --- the entry role belongs to an office competent for the category.
\end{definition}

Restriction (ii) matters for deployment and is cheap to check: a statistically attractive flow
that no officer is authorised to execute is not a policy.

\subsection{The outcome, done carefully}\label{sec:outcome}

This subsection determines what the entire analysis means.

\subsubsection{What the register records}

At closure the officer selects or types a remark. Six standardised remarks account for about
64\% of closures and form a ladder: disposed with no action claimed; disposed with appropriate
action; disposed with the beneficiary benefited. A separate boolean field independently flags
petitioner benefit. The remaining 36\% carry free or semi-standard text.

\subsubsection{Three states, not two}\label{sec:threestate}

The tempting definition is binary: a disposal is \emph{correct} if action was claimed or the
petitioner benefited, incorrect otherwise. This is wrong, in a direction that matters.

Consider the most frequent non-standard closing remarks.

\begin{table}[h]
\centering
\caption{Non-standard closing remarks occurring at least 1,000 times. Aggregate counts only.}
\label{tab:offladder}
\begin{tabular}{lr@{\qquad}lr}
\toprule
Remark & Cases & Remark & Cases\\
\midrule
as reported & 90,061 & not within purview of this cell & 8,344\\
complaint details inadequate & 39,823 & needs a government policy decision & 8,337\\
documents not attached & 28,851 & advised to attend joint hearing & 10,999\\
already taken up / taken up earlier & 21,025 & thanks for the suggestions & 4,769\\
no specific grievance & 16,196 & address not given & 4,086\\
duplicate copy & 15,928 & details not legible & 2,110\\
\bottomrule
\end{tabular}
\end{table}

Read the officer's behaviour off that table. A duplicate copy of an existing grievance; a
complaint with no documents attached; a matter outside the cell's jurisdiction; a request
requiring a policy decision the office cannot make. In every one the officer \emph{did the
right thing} by closing the file, and no action was available. Summing the unambiguous
categories gives roughly 144,700 cases --- about 12\% of all resolved grievances --- correctly
closed and scored as failures by the binary definition.

Worse, the mislabelling is not random with respect to duration. Establishing that a case is a
duplicate is fast. The binary definition therefore removes from the ``correct'' group a large
set of cases that are both correctly handled and quick --- precisely the group a speed analysis
most needs.

\begin{definition}[Actionability and action]\label{def:threestate}
For grievance $i$ define
\[
S_i=\one{\text{the grievance is actionable}},\qquad
C_i=\one{\text{substantive action was taken}}.
\]
$S_i=0$ means no action was available to any officer: duplicate, out of jurisdiction,
insufficient information, requires a policy change. $C_i$ is defined only where $S_i=1$. The
three states are $(S=0)$, $(S=1,C=1)$ and $(S=1,C=0)$; the failure the system exists to reduce
is the last.
\end{definition}

\begin{remark*}[Measurement, and what decided it]
Both $S$ and $C$ are read from the closing remark. Table~\ref{tab:offladder}'s templates plus
the six standard remarks and six further high-frequency strings cover 85.1\% of resolved
closures by exact match; the remainder is free text and is left undetermined rather than
guessed.

The assignments were not decided by reading the wording, which is ambiguous exactly where it
matters. They rest on two structural facts measured over all 1{,}209{,}144 resolved grievances:
screened-out cases run 1.01--1.17 handoffs and close in 1--6 days, while handled cases run
2.18--2.90 handoffs and close in 42--138. The separation is total; nothing high-volume falls
between.

That evidence settles ``as reported'', the most common non-standard remark at 90{,}061 cases and
present in six of the seven intake offices. The closure dropdown photographed in the field record
lists a longer benefit-claiming template beside it, which suggested reading the short string as a
truncation of the long one. \emph{The corpus refutes that}: both longer strings occur fewer than
1{,}000 times and are effectively absent, so there is no co-occurrence to argue from. What the
profile does show is that the case was worked --- 2.18 handoffs and a 78-day median, the longest
of any large template --- while the remark declines to say whether action followed. It is
therefore $S=1$ with $C$ undetermined, a fourth bucket the binary label has nowhere to put, and
it accounts for 7.45\% of resolved grievances.
\end{remark*}

\subsubsection{Why \texorpdfstring{$S$}{S} may be conditioned on and \texorpdfstring{$C$}{C} may not}\label{sec:S-pretreatment}

This is the crux, and it justifies the estimand of \S\ref{sec:estimand}.

$S$ is a property of the grievance. Whether a filing duplicates an earlier one, names a specific
complaint, or falls within the cell's jurisdiction does not depend on which officer receives it.
Formally $S_i(a)=S_i$ for all $a$: the potential values of $S$ do not vary with treatment. A
variable with no treatment effect is, for identification purposes, a pre-treatment covariate,
and conditioning on it is harmless.

$C$ is different. Whether substantive action gets taken is precisely what routing is supposed to
influence, so $C_i(a)$ genuinely varies in $a$. Conditioning on $C=1$ conditions on a
post-treatment variable, and the subpopulation $\{i:C_i(a)=1\}$ is a different set of grievances
for different $a$ --- a \emph{principal stratum}. Comparing mean duration across flows
\emph{within} that set compares different populations, and the comparison has no causal
interpretation even under perfect randomisation.

\begin{example}[Why stratum conditioning fails]\label{ex:stratum}
Suppose flow $a$ takes action on every case, quick and slow alike, while flow $a'$ takes action
only on the easy cases and abandons the hard ones. Then $\E[T(a')\mid C(a')=1]$ is computed over
easy cases only and can be far below $\E[T(a)\mid C(a)=1]$, even though $a$ is better on every
case individually. The comparison rewards $a'$ for shirking.
\end{example}

The design below therefore conditions on $S$ and treats $C$ as an outcome to be constrained,
never as a conditioning event.

\subsection{Duration and censoring}\label{sec:duration}

Let $T_i\ge0$ be days from creation to closure. Data are extracted at a snapshot date, so
grievances still open then have unobserved $T_i$. Let $D_i$ be the censoring time --- days from
arrival to snapshot. Censoring here is \emph{administrative}: $D_i$ is a deterministic function
of the arrival date and the snapshot date. That is a favourable structure, because it makes
Assumption~\ref{ass:censor} credible rather than hopeful.

Censoring is severe in recent cohorts:

\begin{center}
\begin{tabular}{lrrr}
\toprule
Creation year & Grievances & Censored & Rate\\
\midrule
2021 & 51,254 & 1,060 & 2.1\%\\
2022 & 136,902 & 3,098 & 2.3\%\\
2023 & 264,666 & 9,842 & 3.7\%\\
2024 & 667,286 & 61,692 & 9.2\%\\
2025 & 251,180 & 86,452 & 34.4\%\\
\bottomrule
\end{tabular}
\end{center}

A completers-only analysis of the 2025 cohort admits a grievance only if it closed within a few
months, so the analysed subsample is selected on speed and the estimated mean duration is biased
downward by construction.

\begin{definition}[Restricted duration]\label{def:rmst}
Fix a horizon $L>0$ (we use $L=365$ days). The restricted duration is $Y_i=\min(T_i,L)$; the
target is the restricted mean $\E[Y]$.
\end{definition}

Three reasons this is right. It is identified from data on $[0,L]$ alone with no parametric tail
assumption. It is measured in days, so ``the policy saves $\Delta$ days per grievance'' is
literally true. And it equals the area under the survival curve up to $L$
(Lemma~\ref{lem:rmst-area}), which supplies both a nonparametric estimator and a clear picture
of the quantity.

%=====================================================================
\section{The estimand}\label{sec:estimand}
%=====================================================================

\subsection{Potential outcomes and the two functionals}

For each admissible flow $a$, imagine the values realised were grievance $i$ routed through $a$:
$Y_i(a)=\min(T_i(a),L)\in[0,L]$ and $C_i(a)\in\{0,1\}$. We observe the realised flow $A_i$ and,
subject to censoring, the outcomes at that flow.

Restrict attention to actionable grievances, $S=1$. For a flow $a$ and covariate value $x$
define
\begin{equation}
m_a(x)\equiv\E\bigl[Y(a)\mid X=x,\,S=1\bigr],
\qquad
p_a(x)\equiv\Prob\bigl(C(a)=1\mid X=x,\,S=1\bigr).
\label{eq:nuisance}
\end{equation}
In words: $m_a(x)$ is how long an actionable grievance like $x$ would take if routed through
$a$, counting \emph{all} such grievances and not only those that end well; $p_a(x)$ is how
likely substantive action is. Both are ordinary conditional expectations over a fixed
population, with no post-treatment conditioning anywhere.

\subsection{Target and nuisance}\label{sec:nuisance-def}

The functions in \eqref{eq:nuisance}, together with the propensity $e_a(x)$ and censoring
curve $G(t\mid x,a)$ of \S\ref{sec:ident}, are \emph{nuisance functions}: infinite-dimensional
objects that appear in the identification formula for the target but are not themselves of
interest. Nobody wants to know $e_a(x)$, the probability that history would have sent a
grievance like $x$ down flow $a$; that is a description of past bureaucratic habit, not a
policy question. It enters only because the target cannot be written down without it.

The distinction is not bookkeeping --- it determines how accurately each object must be
estimated, and hence which statistical models are admissible for each
(\S\ref{sec:models}). Appendix~\ref{app:semipar} sets out the vocabulary: parametric,
nonparametric and semiparametric models, parameters as functionals, influence functions, the
efficiency bound, and why a target may be estimated at $\sqrt n$ while the nuisances it
depends on are not.

\subsection{Decision rules and their values}

\begin{definition}[Decision rule]\label{def:rule}
A decision rule is a map $\delta:\mathcal X\to\mathcal A$ with $\delta(x)\in\mathcal A(x)$. Write
$\mathcal D$ for the set of such rules.
\end{definition}

$\delta$ is a \emph{function}, not an action: it is the routing protocol, and it is what gets
estimated, evaluated and deployed. The action for a particular grievance is $\delta(x)$. This is
why the value functionals take $\delta$ as their argument.

\begin{definition}[Policy values]\label{def:values}
$V_T(\delta)=\E[Y(\delta(X))\mid S=1]$ and $V_C(\delta)=\E[C(\delta(X))\mid S=1]$.
\end{definition}

$V_T$ is average restricted duration under the rule; $V_C$ is the fraction of actionable
grievances on which action is taken. The historical rule $\delta_{\text{hist}}$ is what the
bureaucracy actually did, and its values are the observed population means.

\subsection{The optimisation problem}

\begin{definition}[Constrained optimal rule]\label{def:optimal}
\begin{equation}
\delta^\star=\argmin_{\delta\in\mathcal D}V_T(\delta)
\qquad\text{subject to}\qquad
V_C(\delta)\ \ge\ V_C(\delta_{\text{hist}}).
\label{eq:program}
\end{equation}
\end{definition}

Read the constraint as the regulator's instruction: \emph{you may reorganise routing to go
faster, but not at the price of taking action less often.} The estimand is the value of solving
the program,
\begin{equation}
\Delta\equiv V_T(\delta_{\text{hist}})-V_T(\delta^\star),
\label{eq:Delta}
\end{equation}
in days saved per actionable grievance, always reported with the realised slack
$V_C(\delta^\star)-V_C(\delta_{\text{hist}})\ge0$.

\subsection{Reduction to a pointwise rule}

Optimising over a function space sounds daunting; it is not, because both functionals are
expectations of pointwise quantities. Define, for $\tau\in[0,1]$,
\begin{equation}
\delta_\tau(x)=\argmin_{a\in\mathcal A(x):\,p_a(x)\ge\tau}m_a(x),
\label{eq:pointwise}
\end{equation}
with the convention that if no $a$ meets the constraint then
$\delta_\tau(x)=\argmax_{a\in\mathcal A(x)}p_a(x)$.

\begin{proposition}[Both values rise with $\tau$]\label{prop:mono}
$V_C(\delta_\tau)$ and $V_T(\delta_\tau)$ are both nondecreasing in $\tau$.
\end{proposition}

\begin{proof}
Appendix~\ref{app:mono}.
\end{proof}

\begin{corollary}[The calibrated floor]\label{cor:feasible}
Let $\mathcal T$ be a finite grid containing $0$ and define
\begin{equation}
\tau^\star=\min\bigl\{\tau\in\mathcal T:\ V_C(\delta_\tau)\ \ge\ V_C(\delta_{\text{hist}})\bigr\},
\label{eq:taustar}
\end{equation}
whenever that set is nonempty. Then $\delta_{\tau^\star}$ is feasible for \eqref{eq:program} and
attains the smallest $V_T$ among feasible rules on the grid.
\end{corollary}

\begin{proof}
Appendix~\ref{app:mono}.
\end{proof}

\begin{remark*}
Proposition~\ref{prop:mono} is the whole economics of the constraint: raising the correctness
floor buys correctness and costs speed, so the calibrated $\tau^\star$ sits exactly where the
constraint first binds. A \emph{larger} $\tau$ than \eqref{eq:taustar} is not conservative ---
it forfeits speed for correctness the regulator did not ask for --- and a smaller one is
infeasible. Fixing $\tau$ at a conventional value, or at zero, and reporting $\Delta$ as though
the constraint had been imposed is a common and consequential error: at $\tau=0$ the constraint
set is everything and \eqref{eq:pointwise} collapses to unconstrained duration minimisation, the
gameable objective \S\ref{sec:intro} warned about.
\end{remark*}

%=====================================================================
\section{Identification}\label{sec:ident}
%=====================================================================

The estimand is written in potential outcomes, which are not observed. The general theory that
recovers such quantities from observational data is in Appendix~\ref{app:causal}; this section
states the assumptions it needs, argues for each in this setting, and assembles the estimating
equation.

\subsection{The assumptions, and what they mean here}

\begin{assumption}[Consistency and no interference]\label{ass:consistency}
If $A_i=a$ then $Y_i=Y_i(a)$ and $C_i=C_i(a)$; and $Y_i(a)$ does not depend on the treatments
assigned to other units.
\end{assumption}

The first clause is innocuous --- it says a flow means the same thing whoever receives it, which
the role-template construction of Definition~\ref{def:flow} is designed to make true. The second
is substantively \emph{false} in a queueing system, and \S\ref{sec:interfere} treats it at
length rather than assuming it away.

\begin{assumption}[Unconfoundedness given the officer's screen]\label{ass:unconf}
For all $a\in\mathcal A(x)$:\quad $\bigl(Y(a),C(a)\bigr)\indep A\mid X,\ S=1$.
\end{assumption}

This is the assumption carrying the weight, and it is only as good as $X$. The covariate list of
\S\ref{sec:covariates} was constructed to reproduce the officer's screen, so that what remains
unobserved is narrow and nameable: routing-private discretion at $t_0$ (``I know this
applicant''), and post-routing shocks (staff leave, a protest, a data-entry backlog). What $U$
is \emph{not} is omitted case difficulty --- difficulty is on the form, and the form is in $X$.
That is a materially better position than the typical routing study, and it is what makes the
sensitivity analysis of \S\ref{sec:sens} interpretable rather than ritual.

The claim is only as strong as the implementation, however. Congestion is on the screen and is
currently absent from the fitted $X$ (\S\ref{sec:backlog}), so any present estimate invokes
Assumption~\ref{ass:unconf} on a strictly smaller information set than the argument for it
describes.

\begin{assumption}[Positivity]\label{ass:positivity}
There is $\epsilon>0$ with $e_a(x)\equiv\Prob(A=a\mid X=x,S=1)\ge\epsilon$ for all $x$ and all
$a\in\mathcal A(x)$.
\end{assumption}

This is the binding constraint in practice. Positivity fails at the flow level in thin cells,
which is what the support floor of Definition~\ref{def:admissible}(i) and the effective sample
size diagnostic of \S\ref{sec:diag} are for. When it fails, the response is to coarsen the
treatment --- back off from flow to department --- not to extrapolate.

\begin{assumption}[Independent censoring]\label{ass:censor}
$D\indep T\mid X,A,S=1$, and $\Prob(D>L\mid X,A,S=1)\ge\epsilon_D>0$.
\end{assumption}

Here administrative censoring earns its keep: $D$ is determined by the arrival date, which is in
$X$, so the conditional independence is close to definitional. The second clause requires that
grievances of every type have had a full $L$ days of exposure --- which is why one does not
evaluate on the newest cohort.

\subsection{What the assumptions buy}

\begin{corollary}[Identification of the estimand]\label{cor:identified}
Under Assumptions~\ref{ass:consistency}--\ref{ass:censor}, the nuisance functions
\eqref{eq:nuisance} satisfy
\[
m_a(x)=\E[Y\mid X=x,A=a,S=1],
\qquad
p_a(x)=\Prob(C=1\mid X=x,A=a,S=1),
\]
the policy values satisfy $V_T(\delta)=\E[m_{\delta(X)}(X)\mid S=1]$ and
$V_C(\delta)=\E[p_{\delta(X)}(X)\mid S=1]$, and consequently $\Delta$ in \eqref{eq:Delta} is a
functional of the observed-data distribution.
\end{corollary}

\begin{proof}
The first display is Theorem~\ref{thm:adjust} applied to $Y$ and to $C$; the second is
Corollary~\ref{cor:value}; censoring is handled by Theorem~\ref{thm:ipcw}. All are proved in
Appendices~\ref{app:causal}--\ref{app:censor}.
\end{proof}

Corollary~\ref{cor:identified} gives \emph{one} route to the estimand, through the outcome
model. Appendix~\ref{app:causal} develops two more --- inverse-probability weighting
(Theorem~\ref{thm:ipw}) and the augmented combination (Theorem~\ref{thm:dr}) --- with different
failure modes. The estimator of \S\ref{sec:est} uses the third.

\subsection{The estimating equation}

Stacking the augmented representation with the censoring weights gives the score used
throughout:
\begin{equation}
\boxed{\;
\psi_i(\delta)
=\hat m_{\delta(X_i)}(X_i)
+\frac{\one{A_i=\delta(X_i)}\ R_i}
      {\hat e_{\delta(X_i)}(X_i)\ \hat G\bigl(Y_i\mid X_i,A_i\bigr)}
\Bigl(Y_i-\hat m_{\delta(X_i)}(X_i)\Bigr)\;}
\label{eq:score}
\end{equation}
where $R_i=\one{D_i\ge Y_i}$ indicates that the restricted duration was observed. Then
$\hat V_T(\delta)=n_S^{-1}\sum_{i:S_i=1}\psi_i(\delta)$. In words: \emph{start with the model's
prediction for the action the rule would take; on those grievances where history happened to
take that same action, correct the prediction by the observed error, scaled up in inverse
proportion to how often that action was taken and how likely the case was to have been observed
at all.}

By Theorem~\ref{thm:dr} the estimator is consistent if the outcome model is right, or if
\emph{both} the propensity and censoring models are right. Replacing $Y$ by $C$ and dropping the
censoring factor gives the corresponding estimator of $V_C$, which is what
Corollary~\ref{cor:feasible} needs to calibrate $\tau$.

%=====================================================================
\section{Estimation}\label{sec:est}
%=====================================================================

We build the estimator in layers, each motivated by a defect in the previous one.

\subsection{Layer 1: outcome regression}\label{sec:layer1}

Corollary~\ref{cor:value} suggests the plug-in
$\hat V^{\text{DM}}_T(\delta)=n_S^{-1}\sum_{i:S_i=1}\hat m_{\delta(X_i)}(X_i)$.

To evaluate this we must compute $\hat m_a(x)$ at a flow the grievance did \emph{not} receive.
That is a design requirement on the feature encoding, and it is easy to violate: if features are
built per dataset --- for instance by encoding categories as positions in a per-sample level
index --- then $\hat m$ is not a well-defined function of $(x,a)$ and the counterfactual
evaluation is meaningless. The encoding must fix one mapping from raw values to model inputs,
estimated once on the training sample and applied unchanged everywhere, with the flow entering
as an explicit argument that can be varied while $x$ is held fixed.

\emph{What goes wrong with only Layer 1.} Nothing checks $\hat m$. If the model errs in a way
correlated with which flows $\delta$ selects --- and it will, because $\delta$ selects flows
precisely where $\hat m$ is small, so selection targets the model's negative errors --- the
value estimate inherits the bias with no diagnostic.

\subsection{Layer 2: weighting}

Theorem~\ref{thm:ipw} suggests instead
$\hat V^{\text{IPW}}_T(\delta)=n_S^{-1}\sum_{i:S_i=1}\one{A_i=\delta(X_i)}\hat e_{\delta(X_i)}(X_i)^{-1}Y_i$,
which uses no outcome model. Its defect is variance: the summand is zero except where history
agreed with the policy, and there it is inflated by $1/\hat e$. If a flow is rare, a handful of
grievances carry the entire estimate.

\subsection{Layer 3: augmentation}

Combining them gives \eqref{eq:score}, with the guarantee of Theorem~\ref{thm:dr}. The
augmentation is a \emph{residual}, mean-zero when the model is right, so Layer 3 reduces to
Layer 1 in the favourable case and to a corrected Layer 2 otherwise.

\subsubsection{Hájek normalisation}

The weights $w_i=\one{A_i=\delta(X_i)}R_i/(\hat e\,\hat G)$ do not average to one in finite
samples. Dividing by their mean rather than by $n$ gives the self-normalised form
\begin{equation}
\hat V^{\text{SN}}_T(\delta)
=\frac1{n_S}\sum_{i:S_i=1}\hat m_{\delta(X_i)}(X_i)
+\frac{\sum_{i:S_i=1}w_i\bigl(Y_i-\hat m_{\delta(X_i)}(X_i)\bigr)}{\sum_{i:S_i=1}w_i}.
\label{eq:hajek}
\end{equation}
This is consistent for the same limit, since the weights converge to mean one, but has bounded
influence: one observation with tiny $\hat e$ moves \eqref{eq:hajek} by at most the residual
range, whereas in the unnormalised form it can move the estimate without limit.

\begin{remark*}[What not to do]
It is tempting to control variance by truncating the \emph{score} $\psi_i$ to a plausible range.
That is not variance reduction but bias: the truncation point interacts with the propensity draw
and the estimator converges to nothing in particular. Trim the \emph{weights} --- equivalently
bound $\hat e$ away from zero, an explicit and reportable restriction of the target population
--- or self-normalise. Never truncate the score.
\end{remark*}

\subsection{Layer 4: cross-fitting}\label{sec:crossfit}

Corollary~\ref{cor:rates} promises that slow-converging nuisances are tolerable, but the
argument has a gap: if $\hat m$ is estimated on the same observations at which it is evaluated,
the two are dependent and an empirical-process term need not vanish. Historically this was
handled by requiring the nuisance estimator to lie in a Donsker class --- a complexity
restriction modern ensembles violate. Cross-fitting removes the requirement entirely
(Proposition~\ref{prop:crossfit}): partition into $K$ folds, fit the nuisances outside each fold,
evaluate the score inside it,
\begin{equation}
\hat V_T(\delta)=\frac1K\sum_{k=1}^K\frac1{|I_k|}\sum_{i\in I_k}\psi_i\bigl(\delta;\hat\eta^{(-k)}\bigr).
\label{eq:crossfit}
\end{equation}
It costs a constant factor in computation and nothing in rate, and should be regarded as
mandatory whenever the nuisances are estimated flexibly.

\subsection{Layer 5: which model estimates what}\label{sec:models}

The theory is agnostic about how $\hat m,\hat e,\hat p,\hat G$ are obtained. The choices are not
arbitrary, and the reasons differ by nuisance.

\begin{table}[h]
\centering
\caption{Nuisance functions, model class, and rationale.}
\label{tab:models}
\begin{tabular}{p{2.2cm}p{3.1cm}p{8.5cm}}
\toprule
Function & Estimator & Why\\
\midrule
$m_a(x)$ \newline duration &
Ridge on $\log(1+T)$ \emph{and} gradient boosting; both reported &
Enters the score as the direct term and inside the residual. Misspecification is tolerable by
Theorem~\ref{thm:dr}, but a \emph{smoother} $m$ reduces the burden carried by $e$. The two
classes bracket the bias--variance trade-off: ridge imposes additive linear structure in the log
scale, boosting imposes none. Reporting both is a specification check --- if $\Delta$ moves
materially between them, functional form is doing the work and must be defended.\\
\addlinespace
$p_a(x)$ \newline action taken &
Penalised logistic or gradient-boosted classifier, with isotonic calibration &
Used only through the threshold $p_a(x)\ge\tau$, so what matters is ranking and
\emph{calibration near $\tau$}, not global fit. An uncalibrated classifier makes $\tau$
uninterpretable, hence the isotonic step.\\
\addlinespace
$e_a(x)$ \newline propensity &
Hierarchical multinomial logistic: department, then flow within department &
Appears in a denominator, so tail behaviour dominates. Tree ensembles produce poorly calibrated
extreme probabilities --- exactly the wrong failure mode. A penalised parametric model with
explicit shrinkage toward the marginal is safer, and the hierarchy matches the institutional
structure while pooling thin cells.\\
\addlinespace
$G(t\mid x,a)$ \newline censoring &
Kaplan--Meier within coarse strata, or Cox on $X$ &
Censoring is administrative, so $G$ is nearly a deterministic function of the arrival date.
Little is gained by flexibility and the stratified estimator is transparent.\\
\bottomrule
\end{tabular}
\end{table}

\subsubsection{Where ridge enters}

The duration model in the log scale is
\[
\log(1+T_i)=\beta'\phi(X_i)+\alpha_{d(i)}+\gamma_{r_0(i)}+\lambda\,n_{\text{esc},i}
+\theta\,\frac{Q_{r_0}(t_{0i})}{\text{cap}_{r_0}}+\varepsilon_i,
\]
with department and entry-role fixed effects, chain length, and normalised congestion. Ridge
rather than ordinary least squares because the design carries many high-cardinality categorical
effects, and the $\ell_2$ penalty shrinks thinly-supported levels toward zero, which is exactly
the partial pooling one wants. The coefficients are interpretable: $\lambda$ is the proportional
cost of an additional handoff, $\gamma_{r_0}$ ranks entry roles by speed. Both are structural
quantities in \eqref{eq:micro}, so the linear specification is not merely convenient --- it is
the estimating equation the mechanism implies.

\subsubsection{Where gradient boosting enters}

Boosting is used for $m$ as the unrestricted comparison. It can represent interactions the
additive form cannot --- that an extra handoff costs more for land disputes than for pension
cases, or that congestion binds only above a threshold. Corollary~\ref{cor:rates} with
cross-fitting makes its use legitimate despite its slow rate.

It must be selected on out-of-sample criteria, and it will not always win. Here the ridge is in
fact the better out-of-sample model (\S\ref{sec:ev-fit}), which is diagnostic rather than
embarrassing: it says flow effects are close to additive in logs, as \eqref{eq:micro} predicts,
and argues for the interpretable model in deployment with boosting retained as the check.

\subsection{Retransformation: the smearing correction}\label{sec:smear}

The duration model is fitted on $\log(1+T)$ but the estimand is on the day scale. Naive
exponentiation is wrong: by Lemma~\ref{lem:jensen-retrans},
$\E[1+T\mid x]=e^{\mu(x)}\E[e^\varepsilon\mid x]\ge e^{\mu(x)}$, so $\exp(\hat\mu(x))-1$
estimates something closer to a conditional \emph{median} than a mean.

\begin{definition}[Smearing estimator]\label{def:smear}
With residuals $\hat\varepsilon_i$ from the log-scale fit, set
$\hat s=n^{-1}\sum_ie^{\hat\varepsilon_i}$ and estimate $\hat m(x)=\hat s\,e^{\hat\mu(x)}-1$.
\end{definition}

$\hat s$ requires no distributional assumption beyond homoskedasticity on the log scale; under
heteroskedasticity compute it within strata. In this application uncorrected predictions
understate the mean by roughly 25--30 days (\S\ref{sec:ev-provisional}), so the correction is not
optional --- and an uncorrected direct term inside \eqref{eq:score} forces the augmentation to
absorb the entire retransformation gap, inflating its variance for no reason.

\subsection{Variance and inference}\label{sec:variance}

Since \eqref{eq:score} is (approximately) an influence function, the asymptotic variance of
$\hat V_T(\delta)$ is $\Var(\psi(O))/n$, estimated by the sample variance of the $\psi_i$ ---
valid under i.i.d.\ sampling, which does not hold here.

Grievances in the same district and year share administrative shocks: a staff vacancy, an
election, a flood. Let $g(i)$ index the district--year cluster; the cluster-robust variance is
\[
\hat\Var\bigl(\hat V_T\bigr)=\frac1{n^2}\sum_g\Bigl(\sum_{i\in g}(\psi_i-\bar\psi)\Bigr)^{\!2},
\]
equivalently obtained by resampling whole clusters. By Lemma~\ref{lem:cluster} the inflation
over the i.i.d.\ variance is approximately $1+(\bar n_g-1)\rho$ with $\rho$ the within-cluster
correlation; with hundreds of grievances per cluster even $\rho=0.01$ doubles the standard
error. Ignoring clustering would materially overstate precision.

%=====================================================================
\section{Learning the rule}\label{sec:policy}
%=====================================================================

So far $\delta$ was fixed and we estimated its value. But $\delta^\star$ must itself be learned,
which introduces a bias easy to miss.

\subsection{The winner's curse}

The empirical rule is $\hat\delta(x)=\argmin_a\hat m_a(x)$ over the constrained set. The
$\argmin$ selects the flow whose \emph{estimate} is smallest, which is systematically the flow
whose estimate is most negatively biased: by Lemma~\ref{lem:winner},
$\E[\min_a\hat m_a]\le\min_a\E[\hat m_a]$. The gap grows with the noise in $\hat m$ and with the
number of candidates, so a policy allowed to choose among many thinly-supported templates will
look better than it is for purely statistical reasons.

\emph{Sample splitting.} Learn $\hat\delta$ on one partition and evaluate $\hat V_T(\hat\delta)$
on another. The evaluation fold sees $\hat\delta$ as a fixed function, so
Corollary~\ref{cor:value} applies unmodified and the winner's curse does not bite. This composes
with cross-fitting: use nested folds, or fit nuisances and policy on complementary folds.

\emph{Restricting the candidate set.} The support floor in Definition~\ref{def:admissible} is
not only about positivity; it also limits the number of candidates and hence the selection bias.
Empirically, restricting to the three most-used admissible flows per cell costs little
achievable value while roughly doubling overlap (\S\ref{sec:ev-provisional}) --- a favourable
trade.

\subsection{Calibrating the constraint}\label{sec:tau}

Corollary~\ref{cor:feasible} requires computing $\tau^\star$ in \eqref{eq:taustar}. Sweep $\tau$
over a grid; for each $\tau$ form $\delta_\tau$ by \eqref{eq:pointwise}; estimate
$V_C(\delta_\tau)$ by the analogue of \eqref{eq:score} with $C$ in place of $Y$; take the
\emph{smallest} $\tau$ meeting the constraint. Report the entire curve
$\tau\mapsto\bigl(\hat V_T(\delta_\tau),\hat V_C(\delta_\tau)\bigr)$: it is the
speed--correctness frontier, far more informative than any single point on it.

\subsection{From a rule to a deployable policy}

Every selected flow must be reachable in the forwarding graph,
Definition~\ref{def:admissible}(ii). And the rule should degrade gracefully, falling back from
flow-level to department-level routing wherever the diagnostics of \S\ref{sec:diag} show
insufficient support.

%=====================================================================
\section{Diagnostics}\label{sec:diag}
%=====================================================================

\subsection{Overlap and effective sample size}

The augmentation in \eqref{eq:score} is a weighted average over the subsample where
$A_i=\delta(X_i)$. Its reliability is governed not by $n$ but by the Kish effective sample size
$\text{ESS}=(\sum_iw_i)^2/\sum_iw_i^2$, which by Lemma~\ref{lem:ess} is the sample size an
unweighted mean would need to match its variance.

Report $\text{ESS}/n$ for every arm. A value of $0.05$ means the correction term carries the
precision of a study one-twentieth the size, and any claim resting on it must be qualified
accordingly. Where $\text{ESS}/n$ falls below a pre-registered threshold, back off to the
coarser department-level treatment, where overlap is far better.

\subsection{Calibration}

$\hat e$ appears in a denominator and $\hat p$ in a threshold, so both need calibration, not
merely discrimination. Report reliability diagrams and the Brier decomposition. A model with
excellent AUC and poor calibration is actively harmful here.

\subsection{Falsification}

Two checks with known answers. \emph{Negative control}: run the pipeline with the outcome
replaced by something routing cannot affect --- the day of the week of filing. A non-zero
estimated effect indicates residual confounding or a leak in feature construction.
\emph{Placebo flows}: permute flow labels within category--district cells and confirm the
estimated $\Delta$ collapses to zero.

%=====================================================================
\section{Sensitivity to unmeasured confounding}\label{sec:sens}
%=====================================================================

Assumption~\ref{ass:unconf} is not testable. The honest response is to quantify how badly it
would have to fail, which the marginal sensitivity model of Appendix~\ref{app:sens} does: a
parameter $\Gamma\ge1$ bounds how much an unmeasured factor may change the odds of receiving a
given flow, \emph{beyond everything on the officer's screen}, with $\Gamma=1$ recovering
unconfoundedness.

For each $\Gamma$ one computes the extremal values of $\hat V_T(\delta)$ over all weight
functions consistent with the model (Proposition~\ref{prop:msm-bounds}). The reportable summary
is the \emph{breakdown point}: the smallest $\Gamma$ at which the interval for $\Delta$ first
contains zero. This converts an untestable assumption into a claim a domain expert can argue
with: \emph{name a case characteristic, visible to the officer but absent from
\S\ref{sec:covariates}, that would shift the odds of the faster routing by a factor of
$\Gamma$.} Because $X$ is designed to be the officer's screen, a large breakdown point is
genuinely reassuring and a small one genuinely alarming.

%=====================================================================
\section{Interference: the queue}\label{sec:interfere}
%=====================================================================

Assumption~\ref{ass:consistency} asserted no interference, which in a queueing system is false.
If the policy sends many grievances to one popular Secretary, that Secretary's backlog rises and
every case routed there --- including cases the policy did not move --- slows down. As defined,
$V_T(\delta)$ is then a partial-equilibrium object: it answers ``what if this one grievance had
been routed differently'', not ``what if the whole system adopted this rule''.

\subsection{A structural model of duration}\label{sec:micro}

The interference has a specific shape, which suggests how to price it. Decompose
\begin{equation}
T_i=\underbrace{\text{wait}\!\left(\frac{Q_{r_0}(t_{0i})}{\text{cap}_{r_0}}\right)}_{\text{queueing delay}}
+\underbrace{\sum_{j=1}^{n_{\text{esc}}}h\bigl(r_j,\text{match}_{ij}\bigr)}_{\text{handling}}
+\underbrace{(n_{\text{esc}}-1)\,h_0}_{\text{handoff cost}}
+\ \varepsilon_i .
\label{eq:micro}
\end{equation}
Only the first term carries interference, and it does so through a scalar state variable
$Q_r(t)$ per role --- the quantity \S\ref{sec:backlog} constructs. Note that \eqref{eq:micro} is
also the justification for the linear specification of \S\ref{sec:models}: the additive
structure in levels motivates an additive approximation in logs, with $\lambda$ estimating $h_0$
in proportional terms.

\subsection{Equilibrium evaluation by replay}

Simulate the arrival stream chronologically. For each grievance draw a counterfactual duration
$Y_i(\delta(X_i))=\hat m_{\delta(X_i)}(X_i)+U_i$ with $U_i$ resampled from the fitted residuals
within a category--district stratum; route it by $\delta$; update $Q_r(t)$ for the entry role,
releasing capacity as cases complete. The resulting $\Delta_{\text{sim}}$ accounts for
congestion and will be smaller than the partial-equilibrium $\Delta$ whenever $\delta$
concentrates load. The gap between them \emph{is} the interference correction, and reporting
both is more informative than either.

A practical variant imposes the load constraint directly in \eqref{eq:pointwise}: forbid the
policy from raising any role's utilisation by more than a set percentage on any day. This
sacrifices some value for a guarantee that the estimated gain does not rest on congestion
effects the model may understate.

%=====================================================================
\section{Evidence}\label{sec:evidence}
%=====================================================================

This section reports what has been computed, organised by how much each result depends on the
design decisions above. Sections~\ref{sec:ev-measure}--\ref{sec:ev-fit} report findings that
stand on their own; \S\ref{sec:ev-provisional} reports policy estimates that do \emph{not} yet
estimate the estimand, and says why.

\subsection{The corpus}

\begin{table}[h]
\centering
\caption{Corpus and chronological splits. ``$C=1$ (binary)'' applies the two-state label of
\S\ref{sec:threestate} that the design replaces; it is shown because every model below was
fitted on it.}
\label{tab:corpus}
\begin{tabular}{lrrrr}
\toprule
Split & Grievances & Resolved & $C=1$ (binary) & Censored\\
\midrule
Train 2021--23 & 452,822 & 438,822 & 162,639 & 3.1\%\\
Validation 2024 & 667,286 & 605,594 & 166,628 & 9.2\%\\
Test 2025 & 251,180 & 164,728 & 34,328 & 34.4\%\\
\midrule
Total & 1,371,288 & 1,209,144 & 363,595 & ---\\
\bottomrule
\end{tabular}
\end{table}

\subsection{Measurement}\label{sec:ev-measure}

These are counting exercises over the register. They depend on no modelling choice, and each is
load-bearing for a design decision made earlier.

\paragraph{The treatment is recoverable (\S\ref{sec:flow}).} Every one of the 1,344,908
escalation chains present in the register decodes to a role sequence; a further 26,380
grievances carry no chain at all, an absent treatment rather than a decoding failure. The chains
use 2,747 distinct user identifiers and collapse to 1,318 role templates over 1,047
category--district cells. After the support floor, 676 cells carry an admissible set, drawing on
111 templates when restricted to the three most common per cell and 266 unrestricted.

\begin{table}[h]
\centering
\caption{Flow census. $n_{\text{esc}}$ is the number of officers who held the file.}
\label{tab:census}
\begin{tabular}{lr@{\qquad\qquad}lr}
\toprule
$n_{\text{esc}}$ & Chains & Most common templates & Chains\\
\midrule
1 & 232,676 & \code{1,5} BDO$\to$Collector & 418,839\\
2 & 664,144 & \code{6} Secretary & 91,838\\
3 & 171,279 & \code{5} Collector & 87,922\\
4 & 242,601 & \code{1,5,6,8} BDO$\to$Collector$\to$Secretary$\to$CMO & 81,195\\
5 & 31,667 & \code{1,37,5,6} & 59,829\\
6--8 & 2,541 & \code{28,5} & 49,126\\
\bottomrule
\end{tabular}
\end{table}

\paragraph{The intake label is not the treatment.} If the citizen's coarse office choice
determined handling there would be no routing decision to optimise.

\begin{center}
\begin{tabular}{lrr}
\toprule
Intake office & Grievances & Ever pending with a Collector\\
\midrule
Collector & 693,691 & 12.8\%\\
Departments & 280,305 & 3.9\%\\
Office of Chief Minister & 217,253 & 49.9\%\\
Chief Secretary & 33,037 & 32.7\%\\
\bottomrule
\end{tabular}
\end{center}

\paragraph{The binary label is mismeasured (\S\ref{sec:threestate}).}
Table~\ref{tab:ladder-evidence} is the evidence for the three-state design. The two rungs the
binary label scores as failures are together 74.8\% of closures, and Table~\ref{tab:offladder}
showed a large part of the second consists of correct closures of non-actionable cases. The
benefit flag also cuts across the rungs rather than agreeing with them: 9.2\% of no-action
closures carry it, while only 70.5\% of closures on the benefit rung do. The binary label is
already a blend of two imperfectly agreeing signals.

\begin{table}[h]
\centering
\caption{Disposal ladder over 1,209,144 resolved grievances.}
\label{tab:ladder-evidence}
\begin{tabular}{lrrrrr}
\toprule
Closing rung & Count & Share & Median $T$ & Mean $T$ & Flagged benefited\\
\midrule
no action claimed & 472,782 & 39.1\% & 46 & 93.9 & 9.2\%\\
non-standard text & 432,222 & 35.7\% & 25 & 71.1 & 3.7\%\\
with appropriate action & 280,887 & 23.2\% & 54 & 96.7 & 6.2\%\\
beneficiary benefited & 23,253 & 1.9\% & 44 & 98.7 & 70.5\%\\
\midrule
Binary $C=1$ & 363,595 & 30.1\% & 54 & 99.5 & ---\\
Binary $C=0$ & 845,549 & 69.9\% & 36 & 80.9 & ---\\
\bottomrule
\end{tabular}
\end{table}

Note the last two rows. Under the binary label, ``correct'' disposals are \emph{slower} than
``incorrect'' ones by 18 days at the median. That is the entire problem in one line: the fastest
way to close a grievance is not to work on it, so an unconstrained speed objective would learn
exactly the wrong policy, and the correctness constraint of \eqref{eq:program} is not a
formality.

\paragraph{Handoffs cost time (\S\ref{sec:micro}).} Table~\ref{tab:nesc-evidence} supports the
$(n_{\text{esc}}-1)h_0$ term in \eqref{eq:micro}: the median roughly doubles between a two-link
and a four-link chain. Beyond four links the pattern is non-monotone on a few thousand cases and
then a few hundred --- a sample-size artefact rather than a reversal.

\begin{table}[h]
\centering
\caption{Duration by chain length, within binary $C=1$.}
\label{tab:nesc-evidence}
\begin{tabular}{lrrrrrrr}
\toprule
$n_{\text{esc}}$ & 1 & 2 & 3 & 4 & 5 & 6 & 7\\
\midrule
Cases & 12,321 & 224,452 & 48,791 & 71,471 & 6,167 & 385 & 8\\
Median $T$ & 37 & 42 & 68 & 98 & 76 & 83 & 133\\
Mean $T$ & 87.6 & 83.8 & 123.3 & 131.6 & 129.7 & 122.7 & 186.4\\
\bottomrule
\end{tabular}
\end{table}

This table also illustrates the identification problem the design must handle and cannot settle
on its own: harder cases are escalated further, so part of the gradient is selection rather than
mechanism. Separating them is what conditioning on $X$ is for, and what \S\ref{sec:sens} probes.

\subsection{Model fit}\label{sec:ev-fit}

\begin{table}[h]
\centering
\caption{Root mean squared error on $\log(1+Y)$, fitted on the actionable population
($S=1$) with the restricted outcome and censoring weights.}
\label{tab:fit}
\begin{tabular}{lccc}
\toprule
& Train 2021--23 & Validation 2024 & Test 2025\\
\midrule
$\hat m$ ridge & 1.174 & \textbf{1.183} & 1.288\\
$\hat m$ gradient boosting & 0.926 & 1.234 & 1.345\\
$\hat m$ boosting, flow features removed & 0.979 & 1.233 & 1.309\\
$\hat m$ ridge, flow features removed & --- & 1.209 & ---\\
\midrule
$\hat p$ AUC, uncalibrated & 0.924 & 0.775 & 0.764\\
$\hat p$ AUC, isotonic & 0.918 & 0.782 & 0.760\\
$\hat p$ ECE, uncalibrated & 0.040 & 0.122 & 0.110\\
$\hat p$ ECE, isotonic & 0.130 & --- & \textbf{0.070}\\
\bottomrule
\end{tabular}
\end{table}

Fitted quantities are not comparable with the previous version of this table, which used the
binary $C=1$ population and an unweighted completers-only target. The isotonic calibration is
fitted on validation, so its validation ECE is in-sample and reported as ``---''; the test
column is the out-of-sample figure.

\emph{The flexible model loses out of sample.} Boosting beats ridge by 0.25 in sample and loses
by 0.05 on validation. This is the empirical content of the claim in \S\ref{sec:models} that
flow effects are close to additive in the log scale, as \eqref{eq:micro} predicts. It argues for
the interpretable model in deployment, and is a reminder that ``use the most flexible learner''
is not an argument --- Corollary~\ref{cor:rates} requires accuracy, not flexibility.

\emph{Whether flow carries predictive signal depends on the learner, and the earlier answer to
this question was contaminated.} An earlier version of this section reported that removing the
flow features cost $0.038$ on validation and called it ``real, modest''. Both the number and
the reading need replacing.

Under the corrected design the boosting ablation is $+0.0003$ --- nothing --- while the ridge
ablation is $+0.0260$ against a cluster-bootstrap standard error of $0.0036$, or seven standard
errors. The asymmetry is interpretable rather than contradictory: a boosted tree can reconstruct
the flow nonparametrically from district, block and category, so the explicit flow columns are
redundant to it; an additive model in the encoded features cannot, so the same columns carry
real information. Any claim that ``the flow does not predict duration'' is a claim about a
learner and must not be restated as a claim about routing.

The $0.038$ itself was largely an artefact. Table~\ref{tab:ablation} decomposes it.

\begin{table}[h]
\centering
\caption{Where the flow-ablation gap comes from. $\Delta$ is
RMSE$_{\text{no flow}}-$RMSE$_{\text{flow}}$ on validation; each rung changes one thing.
Standard errors resample whole district--year clusters.}
\label{tab:ablation}
\begin{tabular}{llccc}
\toprule
Population & Target & $\Delta$ & (eval SE) & (fit SE)\\
\midrule
binary $C=1$, completers & $\log(1+T)$ & $+0.0373$ & 0.0185 & 0.0059\\
actionable $S=1$, completers & $\log(1+T)$ & $-0.0004$ & 0.0085 & ---\\
actionable $S=1$ & $\log(1+Y)$ & $-0.0004$ & 0.0085 & ---\\
actionable $S=1$, IPCW & $\log(1+Y)$ & $-0.0004$ & 0.0085 & 0.0155\\
\midrule
actionable \emph{and} $C=1$, IPCW & $\log(1+Y)$ & $+0.0538$ & 0.0235 & 0.0180\\
\bottomrule
\end{tabular}
\end{table}

The whole movement is the population. The restricted outcome and the censoring weights change
$\Delta$ not at all, which is expected rather than suspicious: every actionable row is resolved
(\S\ref{sec:status}), so $\min(T,L)$ and the capped duration coincide, and the training split is
3.1\% censored so its weights are almost exactly one. Those corrections move the \emph{evaluated}
restricted mean a great deal and the fitted model not at all.

The last row is the diagnosis. Adding $C=1$ back to the actionable population --- changing
nothing else --- restores the gap and exceeds the original. Conditioning on a variable the flow
affects opens a non-causal path between flow and duration, so within that subsample the flow
genuinely is predictive, through selection rather than mechanism. This is
Example~\ref{ex:stratum} showing up as a measurement, and it is the clearest evidence in this
document that the principal-stratum error is not a theoretical scruple.

\emph{The correctness model is usable and is now calibrated.} Validation AUC $0.775$ is adequate
discrimination, and the train--validation gap ($0.924$ to $0.775$) still indicates overfitting.
Since $\hat p$ enters only through the threshold $p_a(x)\ge\tau$, calibration near $\tau$ is the
property that matters; isotonic recalibration on held-out rows reduces test expected calibration
error from $0.110$ to $0.070$ while leaving discrimination essentially unchanged, as a monotone
transformation must.

\subsection{Policy estimates}\label{sec:ev-provisional}

\begin{table}[h]
\centering
\caption{Policy values on the actionable population, restricted mean at $L=365$ with censoring
weights, ridge $\hat m$ with smearing, $\tau=0$. $\Delta_{\text{DM}}$ is the direct-method
contrast, $\Delta_{\text{DR}}$ the augmented one. ``Eligible'' restricts to the three most
common admissible templates per cell.}
\label{tab:provisional}
\begin{tabular}{lrrrrrr}
\toprule
Candidate set & $V_{\text{DM}}$ & $V_{\text{DR}}$ & $\Delta_{\text{DM}}$ & $\Delta_{\text{DR}}$ & (SE) & ESS/$n$\\
\midrule
\multicolumn{7}{l}{\emph{Validation 2024, $n=454{,}232$, IPCW restricted mean 93.5 days}}\\
historical & 107.35 & 93.54 & --- & --- & --- & ---\\
eligible & 96.77 & 73.46 & 10.57 & 20.09 & 7.01 & 0.080\\
unrestricted & 93.25 & 70.09 & 14.09 & 23.45 & 7.99 & 0.051\\
\bottomrule
\end{tabular}
\end{table}

Four of the six gaps that separated the previous version of this table from \eqref{eq:Delta} are
now closed: the population is $S=1$ rather than the binary $C=1$ subset, the outcome is the
restricted mean with censoring weights rather than completers only, predictions are smeared back
to the day scale, and the nuisances are estimated on 2021--23 and evaluated on 2024, which
supplies the independence cross-fitting exists to provide by a stronger route than fold rotation.

The correction to the censored cohort is visible where it was predicted to be. The realised
restricted mean on validation is 93.5 days against a completers-only 84.5; on the 2025 cohort,
which is 34.4\% censored, it is 49.7 against 37.2. The earlier reading that ``realised mean $T$
is 40.6 days against 93.1 on 2024'' was the selection artefact it was suspected of being.

\paragraph{What this establishes.} $\Delta_{\text{DR}}$ is positive on both candidate sets at
roughly three standard errors, so historical routing is not duration-optimal on the actionable
population. The structural reading survives the corrections: restricting from the unrestricted
to the top-three candidate set costs 3.4 days of estimated gain while raising the match rate
from 0.376 to 0.536, so most of the available improvement is reachable within the commonest
flows. ESS$/n$ between 0.051 and 0.080 confirms that flow-level positivity, not sample size, is
the binding statistical constraint.

\paragraph{What it does not establish, and this is the more important half.} Three things.

\emph{The correctness constraint does not resolve.} The historical rate is 0.3868 over the
381{,}436 rows where $C$ is determined. Against that floor the two estimators disagree
qualitatively: the direct method puts the speed-optimising policy at 0.3530, which fails, and
the augmented estimator at 0.4335, which clears it comfortably. $\tau^\star$ is therefore not
identified by this run, and no claim that the policy preserves correctness may be made from it.
Reporting only the augmented figure would be selecting the estimator that gives the convenient
answer.

\emph{The direct and augmented arms still disagree about duration}, by 9.5 days on the eligible
arm. Smearing was expected to close this gap and narrowed it only slightly, so the residual is
not retransformation. Until it is explained, the two arms bracket the answer rather than
agreeing on it.

\emph{Two gaps remain open.} Congestion and trailing performance are still absent from $X$
(\S\ref{sec:backlog}), so Assumption~\ref{ass:unconf} is invoked on a smaller information set
than the argument for it describes, with unknown direction. And $S$ is read from the closing
remark, so it is observable only at closure: the actionable population is conditioned on
resolution, and a grievance still open past the horizon is excluded even though $\min(T,L)$ is
already known for it exactly. IPCW corrects for differential speed among cases that closed; it
cannot restore cases whose $S$ was never observed.

The honest summary is that the sign is now well supported on a defensible population, the
magnitude is bracketed between 10.6 and 23.5 days rather than estimated, and the constraint that
would make any of it a recommendation is unresolved.

%=====================================================================
\section{Implementation status}\label{sec:status}
%=====================================================================

\begin{table}[h]
\centering
\caption{Design components against what exists in
\code{janasunani/experiments/routing\_outcome/}.}
\label{tab:status}
\begin{tabular}{p{6.0cm}p{7.8cm}}
\toprule
Component & Status\\
\midrule
Flow decoding (Def.~\ref{def:flow}) & Built.\\
Shared feature encoding; counterfactual evaluation of $\hat m_a(x)$ (\S\ref{sec:layer1}) & Built.\\
$\hat m$: ridge and gradient boosting & Built.\\
$\hat p$: action classifier & Built, isotonically calibrated on held-out rows.\\
$\hat e$: propensity & Built as empirical category--district shares only; the hierarchical model of Table~\ref{tab:models} is \notbuilt.\\
Augmented score, Hájek normalisation, ESS, cluster bootstrap & Built.\\
Three-state outcome $S,C$ (Def.~\ref{def:threestate}) & Built, for 85.1\% of resolved closures. See the two rows below.\\
Restricted mean and censoring weights (Def.~\ref{def:rmst}, Thm.~\ref{thm:ipcw}) & Built.\\
Smearing (\S\ref{sec:smear}) & Built, per stratum.\\
Cross-fitting (\S\ref{sec:crossfit}) & Supplied by the chronological split: nuisances are fitted on 2021--23 and evaluated on 2024--25, so evaluation rows are out of period rather than merely out of fold.\\
Sample splitting for the rule (\S\ref{sec:policy}) & Built as an option; the winner's curse is priced rather than assumed away.\\
$\tau$ calibration (Cor.~\ref{cor:feasible}) & Built. \emph{Does not resolve}: the direct and augmented estimators disagree about whether the constraint binds (\S\ref{sec:ev-provisional}).\\
\midrule
$S$ for unresolved grievances & \notbuilt. \emph{Highest priority.} $S$ is read from the closing remark, so the actionable population is conditioned on resolution and a case still open past the horizon is dropped although $\min(T,L)$ is known for it. \S\ref{sec:S-pretreatment} licenses predicting $S$ from pre-treatment covariates; nothing does so.\\
Adjudication of the remaining closing remarks & \notbuilt. 14.9\% of resolved closures leave $S$ undetermined and a further 7.5\% leave $C$ undetermined.\\
Congestion $Q_r(t)$ and trailing performance in $X$ (\S\ref{sec:backlog}) & \notbuilt. Weakens Assumption~\ref{ass:unconf} materially.\\
Hierarchical propensity (Table~\ref{tab:models}) & \notbuilt.\\
Sensitivity analysis (\S\ref{sec:sens}) & \notbuilt.\\
Queue replay (\S\ref{sec:interfere}) & \notbuilt.\\
Negative-control and placebo checks (\S\ref{sec:diag}) & \notbuilt.\\
\bottomrule
\end{tabular}
\end{table}

%=====================================================================
\section{Summary}\label{sec:conc}
%=====================================================================

A grievance is a job arriving at a bureaucracy of heterogeneous servers, and the assignment is
made by an officer whose entire information set is recorded. That last fact is what makes the
problem tractable: unconfoundedness given the officer's screen is a claim about unwritten
judgement, not about unobserved case difficulty --- provided the screen is actually reproduced
in $X$, which \S\ref{sec:backlog} shows is not yet the case.

The design has four moving parts, each chosen to avoid a specific failure:

\begin{enumerate}[leftmargin=*]
\item \textbf{A three-state outcome} (\S\ref{sec:threestate}), because a binary
correct/incorrect split scores the correct closure of a duplicate as a failure, and does so
disproportionately among fast cases.
\item \textbf{A constrained rather than conditional estimand} (\S\ref{sec:estimand}), because
conditioning on correctness conditions on a post-treatment variable and compares different
populations across flows (Example~\ref{ex:stratum}).
\item \textbf{An augmented, censoring-weighted, cross-fitted estimator}
(\S\S\ref{sec:ident}--\ref{sec:est}), because outcome regression alone is unchecked, weighting
alone is unstable, censored durations are selected on speed, and flexible nuisances require
sample splitting to be legitimate.
\item \textbf{Explicit accounting for selection in the rule and for congestion}
(\S\S\ref{sec:policy}, \ref{sec:interfere}), because an $\argmin$ over noisy estimates flatters
itself, and a routing policy changes the queues it routes into.
\end{enumerate}

Given Assumptions~\ref{ass:consistency}--\ref{ass:censor}, $\Delta$ is a functional of the
observed-data distribution (Corollary~\ref{cor:identified}), and Theorem~\ref{thm:dr} says the
estimator tolerates misspecification in either the outcome or the treatment model. What remains
is measurement --- above all the three-state outcome and the backlog --- and the estimation
machinery of Table~\ref{tab:status}.

\appendix
\newpage

%=====================================================================
\section{Prerequisites}\label{app:prereq}
%=====================================================================

Everything in the appendices is proved from the following, which are taken as known. References
are to the companion \probrepo{} --- Y.\ Mohanty, \emph{Basic mathematics for statistics and
economics}, the \code{probability} repository --- which develops them from measure-theoretic
foundations.

\begin{enumerate}[leftmargin=*]
\item \textbf{Measure and integration.} Measure spaces, the Lebesgue integral, monotone and
dominated convergence, and Fubini--Tonelli for product measures. \probrepo{}, \emph{Measures},
\emph{Integration}, \emph{Product measures}.
\item \textbf{$L^p$ spaces and inequalities.} The $L^p$ norms; Hölder's inequality (\probrepo{},
\emph{Spaces of functions}, \code{prop:generalizedHolder}) and its $p=q=2$ case Cauchy--Schwarz
(\code{cor:holdersInequality}, \code{prop:innerCauchySchwarz}); Minkowski
(\code{thm:minkowskiInequality}); completeness of $L^p$ (\code{thm:completenessLp}).
\item \textbf{Jensen's inequality} for a convex function and a probability measure. \probrepo{},
\emph{Spaces of functions}, \code{thm:jensenInequality}. Used in
Lemmas~\ref{lem:jensen-retrans} and \ref{lem:winner}.
\item \textbf{Markov's and Chebyshev's inequalities.} \probrepo{}, \emph{Integration},
\code{prop:markovInequality}. Used in Proposition~\ref{prop:crossfit}.
\item \textbf{$L^2$ as a Hilbert space, orthogonal projection, and Riesz representation.}
\probrepo{}, \emph{Spaces of functions}, \code{cor:L2Hilbert}, \code{thm:projectionThm},
\code{cor:orthProjection}, \code{thm:rieszRep}. This is the foundation for conditional
expectation below, and the whole of \S\ref{app:eif} is an application of it: the tangent space
is a closed subspace, the gradient is a Riesz representer, and the efficient influence function
is an orthogonal projection.
\item \textbf{Absolute continuity and Radon--Nikodym derivatives.} The relation $\nu\ll\mu$;
the Radon--Nikodym theorem; the change-of-measure corollary; and the calculus of the derivatives
--- chain rule, sum rule, inverse rule. \probrepo{}, \emph{Differentiation},
\code{def:absoluteContinuityMeasures}, \code{thm:radonNikodym}, \code{cor:radonNikodymIntegral},
\code{prop:RadonNikodymFacts}. Used throughout \S\ref{app:eif} and restated in
\S\ref{app:notation}.
\item \textbf{Independence and product structure.} \probrepo{}, \emph{Independence}.
\item \textbf{Convergence in distribution, the central limit theorem, and Slutsky's theorem.}
Used in Proposition~\ref{prop:ral-clt}. The corresponding \probrepo{} chapters
(\emph{Laws of large numbers}, \emph{Central limit theorems}) are present but not yet
developed, so these are taken as standard.
\item \textbf{Directional derivatives of a functional.} Used in Definition~\ref{def:pathwise}
and Definition~\ref{def:neyman}, both of which differentiate a real-valued map along a
one-parameter family. \probrepo{}, \emph{Differentiation}, reserves
\S\code{sec:differentiation} for differentiation in normed vector spaces, but that section is
not yet written. Nothing beyond the one-variable derivative is actually needed: every such
derivative below is $\frac{d}{dt}$ of a real function of the real parameter $t$, taken at
$t=0$. The Gateaux language is a name for that operation, not an extension of it.
\end{enumerate}

\subsection{Notation: integrals and likelihood ratios}\label{app:notation}

Two conventions are fixed here. The first is an abuse of notation this document declines to use;
the second is item 6 of the list above, restated in the form in which \S\ref{app:eif} needs it.

\textbf{Integrals against a measure.} Many references write the Lebesgue integral of a function
$f$ against a measure $\mu$ as $\int f\,d\mu$. There is no object $d\mu$ in that expression:
$\int\cdot\,d\mu$ is a single indivisible symbol for the integral with respect to $\mu$, the $d$
inherited by analogy from Riemann's $\int f(x)\,dx$ and carrying no meaning of its own. In
particular it is unrelated to the $d$ of differential forms, which the resemblance invites and
which plays no part anywhere in this document. \probrepo{} avoids it too, writing the integral as
$\bar\mu(f)$ --- the measure applied to the function, which is what it is --- and reserving
$\int f\,d\mu$ for the remark that introduces Radon--Nikodym derivatives, where it is labelled
``traditional notation''.

Here the measure being integrated against is a probability measure at every occurrence, so the
integral is an expectation and the expectation is what is written: $\bar P(f)=\E_P[f(O)]$. The
symbol $\int$ is reserved for ordinary integration against length on the real line, as in
$\int_0^Lg(t)\,dt$ (Lemma~\ref{lem:rmst-area}), where $dt$ does mark a variable of integration.

\textbf{Likelihood ratios.} Write $Q\ll P$ when $Q$ is absolutely continuous with respect to $P$,
meaning $P(B)=0$ forces $Q(B)=0$ (\probrepo{}, \emph{Differentiation},
\code{def:absoluteContinuityMeasures}, where the relation is written $Q<<P$). For $\sigma$-finite
$Q,P$ the Radon--Nikodym theorem (\code{thm:radonNikodym}) says $Q\ll P$ holds if and only if
there is a nonnegative measurable function, unique up to $P$-null sets, reproducing $Q$ as an
average against $P$; and its corollary (\code{cor:radonNikodymIntegral}) converts that into the
change-of-measure rule. In the notation of this document, that function --- written $dQ/dP$ and
called the \emph{Radon--Nikodym derivative}, or in this context the \emph{likelihood ratio} of
$Q$ to $P$ --- is characterised by
\begin{equation}
Q(B)=\E_P\left[\one{O\in B}\,\frac{dQ}{dP}(O)\right]\ \text{ for every measurable }B,
\qquad\text{equivalently}\qquad
\E_Q[f]=\E_P\left[f\,\frac{dQ}{dP}\right]
\label{eq:rn}
\end{equation}
for every $f$ integrable under $Q$. The second identity is the only use made of densities below.

Two remarks, both from \probrepo{}. First, $dQ/dP$ is a derivative only in a stretched sense ---
it is not defined as a limit of ratios except against Lebesgue measure --- and what earns it the
Leibniz notation is that it obeys the expected calculus (\code{prop:RadonNikodymFacts}): a chain
rule $\frac{dR}{dP}=\frac{dR}{dQ}\frac{dQ}{dP}$ when $R\ll Q\ll P$, a sum rule, and
$\frac{dP}{dQ}=(\frac{dQ}{dP})^{-1}$ when the two measures are mutually absolutely continuous.
The chain rule is what lets a path be reparameterised in \S\ref{app:eif} without recomputing
anything. Second, the reference measure is written explicitly and is always the fixed $P$ below;
a bare $dP$ never appears.

Because these ratios get differentiated again with respect to a path parameter, and stacking
$\partial/\partial t$ on top of $dQ/dP$ is unreadable, they are given a letter: $L\equiv dQ/dP$.
Every statement about $L$ below is a statement about the function characterised by
\eqref{eq:rn}.

\subsection{Conditional expectation}\label{app:condexp}

One prerequisite is used constantly below and is not yet developed in \probrepo{}, so it is
stated here. For $Y\in L^1(\Omega,\mathcal F,\Prob)$ and a sub-$\sigma$-algebra
$\mathcal G\subseteq\mathcal F$, the conditional expectation $\E[Y\mid\mathcal G]$ is the
$\Prob$-a.s.\ unique $\mathcal G$-measurable integrable random variable satisfying
$\E\bigl[\one{G}\,\E[Y\mid\mathcal G]\bigr]=\E\bigl[\one{G}\,Y\bigr]$ for every $G\in\mathcal G$
--- averaging the conditional expectation over any event that $\mathcal G$ can resolve gives the
same answer as averaging $Y$ itself. Existence
follows from the Radon--Nikodym theorem, or for $Y\in L^2$ from orthogonal projection onto the
closed subspace of $\mathcal G$-measurable square-integrable functions (item 5 above). Three
properties are used without further comment:

\begin{enumerate}[leftmargin=*]
\item \textbf{Tower property.} For $\mathcal H\subseteq\mathcal G$,
$\E\bigl[\E[Y\mid\mathcal G]\bigm|\mathcal H\bigr]=\E[Y\mid\mathcal H]$; in particular
$\E\bigl[\E[Y\mid X]\bigr]=\E[Y]$.
\item \textbf{Taking out what is known.} If $Z$ is $\mathcal G$-measurable and $ZY$ integrable,
$\E[ZY\mid\mathcal G]=Z\,\E[Y\mid\mathcal G]$.
\item \textbf{Independence.} If $Y\indep\mathcal G$ then $\E[Y\mid\mathcal G]=\E[Y]$.
\end{enumerate}

Throughout the appendices we work conditional on $S=1$ and suppress it from the notation.

%=====================================================================
\section{Parametric, nonparametric, and semiparametric models}\label{app:semipar}
%=====================================================================

This appendix fixes the vocabulary used throughout: what a statistical model is, what a
parameter is when the model is not parametric, what makes something a nuisance, and what
``efficient'' means. Nothing here is specific to routing.

\subsection{Models}

\begin{definition}[Statistical model]\label{def:model}
A \emph{statistical model} is a set $\mathcal P$ of candidate probability distributions for
the observed data $O$. The data are assumed drawn from some $P\in\mathcal P$, and the
analyst's assumptions are exactly the restrictions that define $\mathcal P$.
\end{definition}

Models are classified by how they are indexed.

\begin{itemize}[leftmargin=*]
\item \textbf{Parametric.} $\mathcal P=\{P_\theta:\theta\in\Theta\}$ with
$\Theta\subseteq\mathbb R^d$ finite-dimensional. Example: $O\sim\mathcal N(\mu,\sigma^2)$,
$\theta=(\mu,\sigma^2)$. Everything about $P$ is determined by $d$ numbers.
\item \textbf{Nonparametric.} $\mathcal P$ is all distributions on the sample space, subject
only to mild regularity (moments, absolute continuity). No finite-dimensional index exists.
\item \textbf{Semiparametric.} $\mathcal P$ is indexed by a pair $(\theta,\eta)$ with
$\theta\in\mathbb R^d$ finite-dimensional and $\eta$ infinite-dimensional. The classical
example is the Cox model, where the coefficient vector is $\theta$ and the baseline hazard is
$\eta$.
\end{itemize}

The word ``semiparametric'' is also used, more loosely and more usefully for our purposes, for
any problem in which \emph{the model is large but the quantity of interest is
finite-dimensional}. That is our situation exactly: we assume nothing about the shape of the
distribution of $(X,A,Y)$ --- the model is nonparametric --- yet the object we report,
$\Delta$, is a single number. The tension between a rich model and a low-dimensional target is
what the rest of this appendix is about.

\subsection{Parameters as functionals}

In a parametric model, ``the parameter'' is the index $\theta$. In a nonparametric model there
is no index, so the word must be redefined.

\begin{definition}[Parameter, target]\label{def:functional}
A \emph{parameter} of the model $\mathcal P$ is a map $\psi:\mathcal P\to\mathbb R^k$. The
\emph{target} is the parameter one wishes to report.
\end{definition}

Thus $V_T(\delta)=\E_P\bigl[m_{\delta(X)}(X)\bigr]$ is a parameter in this sense: a rule assigning
a number to each candidate distribution. This viewpoint is not pedantry. It makes the question
``how well can this be estimated?'' well posed even when the model has no finite-dimensional
index, because one can ask how much $\psi(P)$ moves as $P$ is perturbed --- which is exactly
what \S\ref{app:eif} does.

\begin{definition}[Nuisance]\label{def:nuisance}
Given a target $\psi$, a \emph{nuisance} is any further component of $P$ required to evaluate
$\psi$ but not itself reported. A \emph{nuisance parameter} is finite-dimensional; a
\emph{nuisance function} is not.
\end{definition}

The vocabulary is inherited from parametric statistics: testing a hypothesis about a normal
mean $\mu$, the variance $\sigma^2$ must be dealt with and is never the object of inference.
Here the same logic applies to whole functions --- $m_a(\cdot)$, $e_a(\cdot)$, $G(\cdot)$.

\subsection{Why the split matters}\label{app:whysplit}

Three consequences run through the whole document, and they are the reason the estimator is
built as it is.

\begin{enumerate}[leftmargin=*]
\item \emph{Accuracy requirements are weaker for nuisances than for the target.} We need the
target estimated consistently at $\sqrt n$ with a valid confidence interval. We do \emph{not}
need the nuisances good in any absolute sense --- only good enough that their errors do not
contaminate the target. Theorem~\ref{thm:dr} makes ``good enough'' precise and the answer is
remarkably weak: the bias is the \emph{product} of two errors, so each may converge at the
slow rate $n^{-1/4}$ (Corollary~\ref{cor:rates}).
\item \emph{Hence flexible machine learning is admissible for nuisances.} A gradient-boosted
ensemble has no $\sqrt n$ theory and no usable standard errors, which would disqualify it if it
were estimating the target. As a nuisance estimator it is fine, provided it is cross-fitted
(Proposition~\ref{prop:crossfit}). This is the entire licence for \S\ref{sec:models}.
\item \emph{Robustness is stated in terms of which nuisance is wrong.} ``Doubly robust'' means
precisely: consistent if the outcome nuisance is right \emph{or} the treatment nuisance is
right. The guarantee is meaningful only because the two are separate objects with separate
failure modes.
\end{enumerate}

\begin{remark*}
A decision rule $\delta$ is emphatically \emph{not} a nuisance. It is the deliverable --- the
protocol that would be adopted --- and it is judged by different criteria: interpretability,
administrative feasibility, stability. None of those matter for $\hat m$ or $\hat e$. Rule of
thumb: if a component would appear in the deployed system it is not a nuisance; if it exists
only inside the estimator, it is.
\end{remark*}

\subsection{Influence functions}\label{app:eif}

Two questions have one answer. \emph{How accurately can $\psi$ be estimated at all, by any
method whatever?} And \emph{what does the error of a particular estimator look like when $n$ is
large?} The object that answers both is the influence function. The term carries two meanings
--- one attached to a functional, one to an estimator --- which are ordinarily introduced
together and left to merge. They are taken separately here and then connected, because the
document uses both.

\paragraph{Differentiating a functional.}
A parameter is a map $\psi:\mathcal P\to\mathbb R$ (Definition~\ref{def:functional}). Asking how
hard it is to estimate is asking how fast $\psi$ moves when $P$ moves: if two distributions are
nearly indistinguishable from $n$ observations but assign very different values to $\psi$, then
no procedure can reliably tell those values apart. That is a statement about a derivative. So we
must differentiate a function whose argument is a probability measure --- and $\mathcal P$ has
no coordinates to differentiate along. The standard device is to restrict to one-dimensional
families through $P$ and differentiate along each one.

\begin{definition}[Path and score]\label{def:path}
A \emph{path} through $P\in\mathcal P$ is a family $\{P_t:|t|<\epsilon\}\subseteq\mathcal P$ with
$P_0=P$ and $P_t\ll P$ for each $t$ --- so that the Radon--Nikodym theorem
(Appendix~\ref{app:prereq}, item 6) supplies the likelihood ratios $L_t\equiv dP_t/dP$
characterised by \eqref{eq:rn} --- such that $L_0=1$ and $L_t$ admits a derivative
\[
h(o)\ \equiv\ \left.\frac{\partial}{\partial t}\log L_t(o)\right|_{t=0}
\]
lying in $L^2(P)$. The function $h$ is the \emph{score} of the path.
\end{definition}

Read $h$ as the direction of travel: it records, for each possible observation $o$, the
proportional rate at which the model reweights $o$ as $t$ moves off zero. Since
$P_t$ is a probability measure, \eqref{eq:rn} with $B$ the whole space gives $\E_P[L_t]=1$ for
every $t$; differentiating at $t=0$ and using $\partial_tL_t=L_t\,\partial_t\log L_t$ gives
$\E_P[h]=0$. Every score therefore lies in
\[
L^2_0(P)\ \equiv\ \bigl\{h\in L^2(P):\E_P[h]=0\bigr\},
\]
a closed subspace of $L^2(P)$ and hence itself a Hilbert space under
$\langle f,g\rangle=\E_P[fg]$ (Appendix~\ref{app:prereq}, item 5). A mean-zero direction is one
that shifts weight around without creating any, which is what keeps total probability at one.

One construction supplies all the paths needed below: for bounded $h$ with $\E_P[h]=0$, take the
\emph{tilted path}
\begin{equation}
L_t(o)=1+t\,h(o),
\qquad\text{that is,}\qquad
P_t(B)=\E_P\bigl[\one{O\in B}\bigl(1+t\,h(O)\bigr)\bigr],
\label{eq:tiltpath}
\end{equation}
which is a probability measure for $|t|<1/\|h\|_\infty$: it is nonnegative by the bound, and has
total mass $\E_P[1+th]=1$ because $\E_P[h]=0$. Its score is $h$, since
$\partial_t\log(1+th)|_{t=0}=h$. Every bounded mean-zero direction is therefore achievable, as
long as the tilted measure stays inside $\mathcal P$. Which directions survive that requirement
is exactly where the model's assumptions show up.

\begin{definition}[Tangent space]\label{def:tangent}
The \emph{tangent space} $\dot{\mathcal P}(P)\subseteq L^2_0(P)$ is the closed linear span of the
scores of all paths through $P$ that remain in $\mathcal P$.
\end{definition}

Two extremes bracket every case.

\begin{itemize}[leftmargin=*]
\item \textbf{Nonparametric model.} $\mathcal P$ is unrestricted, so \eqref{eq:tiltpath} stays in
the model for every bounded mean-zero $h$, and those are dense in $L^2_0(P)$. The tangent space
is all of $L^2_0(P)$: $P$ may be perturbed in every direction.
\item \textbf{Parametric model} $\{P_\theta:\theta\in\Theta\subseteq\mathbb R^d\}$, with
$P=P_{\theta_0}$ and $L_\theta\equiv dP_\theta/dP$. The only paths are
curves $t\mapsto P_{\theta(t)}$ in $\Theta$, whose scores are, by the chain rule, linear
combinations of the $d$ coordinate scores $\partial\log L_\theta/\partial\theta_j$ evaluated at
$\theta_0$ --- the classical score vector of maximum-likelihood theory. The tangent space is
their $d$-dimensional span.
\end{itemize}

The tangent space is the geometry of what has been assumed: the more $\mathcal P$ is restricted,
the fewer directions remain, and the smaller $\dot{\mathcal P}(P)$. Everything about the
assumptions-versus-precision trade-off follows from that one sentence.

\begin{definition}[Pathwise differentiability, gradient]\label{def:pathwise}
$\psi$ is \emph{pathwise differentiable} at $P$ if there exists $\varphi\in L^2_0(P)$ with
\begin{equation}
\left.\frac{d}{dt}\psi(P_t)\right|_{t=0}=\E_P\bigl[\varphi(O)\,h(O)\bigr]
\label{eq:pathwise}
\end{equation}
for every path through $P$ in $\mathcal P$, $h$ being that path's score. Any such $\varphi$ is a
\emph{gradient} of $\psi$ at $P$.
\end{definition}

Unpacked, \eqref{eq:pathwise} says two things. The derivative along a path depends on the path
only through its score --- so the score is a legitimate notion of ``direction''. And the map
$h\mapsto\frac{d}{dt}\psi(P_t)|_{t=0}$ is a bounded linear functional on the tangent space. The
Riesz representation theorem (Appendix~\ref{app:prereq}, item 5) says every such functional on a
Hilbert space is an inner product against a fixed element, and that element is the gradient. So
pathwise differentiability is not an exotic requirement: given that the derivative exists and
depends boundedly on the direction, a gradient exists automatically. This is the sense in which
$\varphi$
\emph{is} the derivative of $\psi$ at $P$: not a number, but the vector representing the
directional derivative in all directions at once.

Nothing so far makes the gradient unique. It is not, and the way it fails is the whole point.

\begin{proposition}[Canonical gradient]\label{prop:canonical}
Let $\psi$ be pathwise differentiable at $P$ with gradient $\varphi$. Then
\begin{enumerate}[leftmargin=*,label=(\roman*)]
\item $\tilde\varphi\in L^2_0(P)$ is a gradient if and only if
$\tilde\varphi-\varphi\perp\dot{\mathcal P}(P)$;
\item exactly one gradient lies in $\dot{\mathcal P}(P)$, namely the orthogonal projection
$\varphi^{\mathrm{eff}}=\Pi\bigl(\varphi\mid\dot{\mathcal P}(P)\bigr)$;
\item $\E_P[(\varphi^{\mathrm{eff}})^2]\le\E_P[\tilde\varphi^2]$ for every gradient
$\tilde\varphi$, with equality only when $\tilde\varphi=\varphi^{\mathrm{eff}}$.
\end{enumerate}
\end{proposition}

\begin{proof}
(i) Two gradients reproduce the same derivative against every score, so
$\E[(\tilde\varphi-\varphi)h]=0$ for all $h$ in a spanning set of $\dot{\mathcal P}(P)$, hence
for all of it by linearity and continuity of the inner product; conversely, adding to $\varphi$
anything orthogonal to the tangent space changes no such inner product, so \eqref{eq:pathwise}
still holds.

(ii) The projection onto a closed subspace exists and is unique (Appendix~\ref{app:prereq}, item
5), and $\varphi-\varphi^{\mathrm{eff}}\perp\dot{\mathcal P}(P)$ by construction, so
$\varphi^{\mathrm{eff}}$ is a gradient by (i). If two gradients both lay in
$\dot{\mathcal P}(P)$, their difference would be both in the tangent space, being a difference of
its elements, and orthogonal to it by (i); the only such element is $0$.

(iii) Write $\tilde\varphi=\varphi^{\mathrm{eff}}+(\tilde\varphi-\varphi^{\mathrm{eff}})$. The
second term is orthogonal to $\dot{\mathcal P}(P)$ by (i) and the first lies in it, so Pythagoras
gives $\|\tilde\varphi\|^2=\|\varphi^{\mathrm{eff}}\|^2+\|\tilde\varphi-\varphi^{\mathrm{eff}}\|^2$,
and the cross term vanishes only when the two coincide.
\end{proof}

\begin{remark*}
In a nonparametric model $\dot{\mathcal P}(P)=L^2_0(P)$, the projection in (ii) is the identity,
and the gradient is unique --- there is nothing to choose. That is our situation
(\S\ref{app:semipar}), so throughout this document ``the'' influence function is well defined
and the qualifier ``efficient'' adds no content. The distinction matters only when a model is
genuinely restricted, which is precisely when a gradient can be improved by projecting away the
part of it the model forbids.
\end{remark*}

\paragraph{Where the name comes from.}
Take the contamination path $P_t=(1-t)P+t\,\delta_o$, where $\delta_o$ places all its mass on
the single point $o$: a fraction $t$ of the population is replaced by copies of one observation.
It is not literally of the form \eqref{eq:tiltpath}. If $P$ has no atom at $o$ then $\delta_o$ is
not merely outside the reach of $P$ but mutually singular with it (\probrepo{},
\emph{Differentiation}, \code{def:mutualSingularity}): the two measures concentrate on disjoint
sets, $\{o\}$ and its complement. So $\delta_o\not\ll P$, no likelihood ratio exists, and
Definition~\ref{def:path} does not apply --- the Lebesgue decomposition of $P_t$ relative to $P$
(\code{thm:lebesgueDecomposition}) has a singular part of size $t$ that never becomes a density.
What survives is the derivative itself: $t\mapsto\psi(P_t)$ is an ordinary real function of a
real variable, and for the functionals here it is differentiable at $t=0$ regardless. Taking that
derivative gives the formula that \eqref{eq:pathwise} would have produced,
\begin{equation}
\varphi(o)=\lim_{t\downarrow0}\frac{\psi\bigl((1-t)P+t\,\delta_o\bigr)-\psi(P)}{t}.
\label{eq:hampel}
\end{equation}
The gradient evaluated at $o$ is the effect on the parameter of contaminating the population with
an infinitesimal mass of observations at $o$: the \emph{influence} of an observation at $o$,
which is where the name and Hampel's original robustness motivation come from. An influence
function with a heavy tail is an estimator that a few observations can move --- exactly what the
overlap diagnostics of \S\ref{sec:diag} watch for, since a small $\hat e$ makes
\eqref{eq:score} large at that point. Equation~\eqref{eq:hampel} is also the practical recipe for
computing one: write $\psi$ as an explicit functional of $P$, substitute the contaminated law,
differentiate at $t=0$.

\begin{example}[The mean]\label{ex:if-mean}
$\psi(P)=\E_P[Y]$. Along \eqref{eq:tiltpath}, the change-of-measure rule \eqref{eq:rn} gives
$\psi(P_t)=\E_{P_t}[Y]=\E_P[Y\,L_t]=\E_P\bigl[Y(1+th)\bigr]=\E_P[Y]+t\,\E_P[Yh]$,
so the derivative at $t=0$ is
$\E_P[Yh]=\E_P\bigl[(Y-\E_P[Y])h\bigr]$, the second equality using $\E_P[h]=0$. Comparing with
\eqref{eq:pathwise}, $\varphi(O)=Y-\E_P[Y]$. The sample mean then satisfies
\eqref{eq:asymplinear} below exactly, with a remainder that is identically zero, and its variance
is $\Var_P(Y)$ --- which by Proposition~\ref{prop:canonical} and the remark above is the bound.
No estimator of a population mean improves on the sample mean when nothing is assumed.
\end{example}

\begin{example}[Policy value]\label{ex:if-value}
For $\psi(P)=V_T(\delta)=\E[Y(\delta(X))]$ in the nonparametric model, identified under
Assumptions~\ref{ass:consistency}--\ref{ass:positivity} and with no censoring, the gradient is
\[
\varphi(O)=\psi\bigl(O;m^\ast,e^\ast,\delta\bigr)-V_T(\delta),
\]
the augmented score \eqref{eq:aipw} centred at the estimand. The computation is
\eqref{eq:hampel} applied to the adjustment formula of Theorem~\ref{thm:adjust}; the inverse
propensity appears because contaminating at $o$ perturbs the conditional law of $A$ given $X$ as
well as that of $Y$ given $(X,A)$, and the augmentation term is what corrects for the second
channel. This is the reason \eqref{eq:score} has the shape it does.
\end{example}

\paragraph{Influence functions of estimators.}
The second sense of the term describes not a functional but a procedure.

\begin{definition}[Asymptotically linear estimator]\label{def:ral}
An estimator $\hat\psi$ of $\psi(P)$ is \emph{asymptotically linear} with \emph{influence
function} $\varphi$ if $\E_P[\varphi(O)]=0$, $\E_P[\varphi(O)^2]<\infty$, and
\begin{equation}
\sqrt n\bigl(\hat\psi-\psi(P)\bigr)=\frac1{\sqrt n}\sum_{i=1}^n\varphi(O_i)+R_n,
\qquad R_n\ \xrightarrow{\ \Prob\ }\ 0 .
\label{eq:asymplinear}
\end{equation}
It is \emph{regular} if its limit distribution is unchanged under local perturbations of $P$
--- a condition that excludes pathological superefficient estimators such as Hodges'.
\end{definition}

The content of \eqref{eq:asymplinear} is that a possibly complicated estimator, built by fitting
models and solving equations, behaves in large samples like a simple average of a fixed function
of one observation at a time. Everything about its first-order behaviour is then read off that
function.

\begin{proposition}\label{prop:ral-clt}
If $\hat\psi$ is asymptotically linear with influence function $\varphi$, then
$\sqrt n(\hat\psi-\psi(P))\rightsquigarrow\mathcal N\bigl(0,\E_P[\varphi(O)^2]\bigr)$.
\end{proposition}

\begin{proof}
The leading term of \eqref{eq:asymplinear} is a normalised sum of i.i.d.\ mean-zero,
finite-variance random variables, so the central limit theorem gives convergence to
$\mathcal N(0,\E[\varphi^2])$; the remainder converges to zero in probability, so Slutsky's
theorem leaves the limit unchanged.
\end{proof}

Proposition~\ref{prop:ral-clt} is why influence functions are the right currency: the
asymptotic variance of an asymptotically linear estimator \emph{is} the second moment of its
influence function. It is also the justification for the variance estimator of
\S\ref{sec:variance} --- once the score \eqref{eq:score} is recognised as an influence
function, the sample variance of the scores estimates the asymptotic variance directly, with
no delta method.

\begin{remark*}[The two senses agree]
If $\hat\psi$ is regular and asymptotically linear with influence function $\varphi$ in the sense
of Definition~\ref{def:ral}, then $\varphi$ is a gradient of $\psi$ at $P$ in the sense of
Definition~\ref{def:pathwise}. This is the one result quoted rather than proved here --- it needs
local asymptotic normality and Le Cam's third lemma, neither of which is developed in
\probrepo{}; see van der Vaart, \emph{Asymptotic Statistics}, Ch.~25. It is what licenses the
single word for the two definitions, and it is what makes a lower bound on the variance of
gradients a lower bound over estimators. Nothing proved below depends on it.
\end{remark*}

\begin{definition}[Efficient influence function, efficiency bound]\label{def:bound}
The \emph{efficient influence function} (EIF) of $\psi$ in $\mathcal P$ at $P$ is the canonical
gradient $\varphi^{\mathrm{eff}}$ of Proposition~\ref{prop:canonical}(ii). Its variance
$\E_P[(\varphi^{\mathrm{eff}})^2]$ is the \emph{semiparametric efficiency bound} for $\psi$ in
$\mathcal P$.
\end{definition}

By the remark above and Proposition~\ref{prop:canonical}(iii), no regular asymptotically linear
estimator has asymptotic variance below the bound: its influence function is a gradient, and
every gradient is at least as variable as the canonical one. This is the infinite-dimensional
analogue of the Cramér--Rao inequality, and the practical consequence is that an estimator built
from the EIF is optimal, so the search for a better one may stop.

\begin{corollary}[Assumptions buy precision]\label{cor:nested}
Let $\mathcal P_1\subseteq\mathcal P_2$ be two models containing $P$, in both of which $\psi$ is
pathwise differentiable at $P$. Then the efficiency bound for $\psi$ in $\mathcal P_1$ is no
larger than the bound in $\mathcal P_2$.
\end{corollary}

\begin{proof}
Fewer paths lie in $\mathcal P_1$ than in $\mathcal P_2$, so
$\dot{\mathcal P}_1(P)\subseteq\dot{\mathcal P}_2(P)$ and \eqref{eq:pathwise} is a weaker
requirement in $\mathcal P_1$: any gradient in $\mathcal P_2$ is a gradient in $\mathcal P_1$. In
particular $\varphi_2^{\mathrm{eff}}$ is a gradient in $\mathcal P_1$, so
Proposition~\ref{prop:canonical}(iii) applied within $\mathcal P_1$ gives
$\E[(\varphi_1^{\mathrm{eff}})^2]\le\E[(\varphi_2^{\mathrm{eff}})^2]$.
\end{proof}

That is the formal statement of the trade-off between assumptions and precision, and it cuts
the way one would hope: assuming more can only tighten the achievable interval, and the
nonparametric bound --- our case --- is the most conservative of all. It is the price paid for
the credibility argued in \S\ref{sec:ident}, and it is a price, not a free choice.

\paragraph{From an influence function to an estimator.}
Gradients are not only a yardstick; they tell one how to build the estimator. The plug-in
$\psi(\hat P)$ formed from a flexible fit is generally not $\sqrt n$-consistent, because the
regularisation that makes $\hat P$ a good density estimate biases the functional. Pathwise
differentiability supplies the correction. For a second distribution $\bar P$, expanding $\psi$
along the path from $\bar P$ to $P$ and collecting the first-order term gives the \emph{von Mises
expansion}
\begin{equation}
\psi(\bar P)-\psi(P)=-\E_P\bigl[\varphi(O;\bar P)\bigr]+R_2(\bar P,P),
\label{eq:vonmises}
\end{equation}
where $\varphi(\cdot;\bar P)$ is the gradient evaluated at $\bar P$ and $R_2$ collects terms of
second order in $\bar P-P$. Rearranging and replacing the unknown $\E_P$ by the sample average
gives the \emph{one-step estimator}
\begin{equation}
\hat\psi=\psi(\hat P)+\frac1n\sum_{i=1}^n\varphi\bigl(O_i;\hat P\bigr):
\label{eq:onestep}
\end{equation}
fit, then add the empirical mean of the estimated influence function. Its error is the average of
$\varphi(O_i;P)$, an empirical-process term, and $R_2$ --- and $R_2$ being second order is what
makes the whole construction work.

The estimator of this document is exactly \eqref{eq:onestep}. Take $\psi(\hat P)$ to be the
outcome-regression plug-in $n^{-1}\sum_i\hat m_{\delta(X_i)}(X_i)$ and $\varphi$ the gradient of
Example~\ref{ex:if-value}; the two plug-in terms cancel and what remains is the average of the
augmented score \eqref{eq:aipw}. The three pieces then line up:

\begin{itemize}[leftmargin=*]
\item $R_2$ in \eqref{eq:vonmises} \emph{is} the product bias \eqref{eq:productbias} of
Theorem~\ref{thm:dr}, which is why that bias is a product of two nuisance errors rather than a
sum;
\item the empirical-process term is what cross-fitting controls
(Proposition~\ref{prop:crossfit});
\item and the statement that $R_2$ has no first-order part is Neyman orthogonality
(\S\ref{app:neyman}).
\end{itemize}

Double robustness, orthogonality, and the second-order remainder in \eqref{eq:vonmises} are one
property described three ways, and the reason to prefer this estimator is that it is built from
the gradient rather than in spite of it.

The score \eqref{eq:score} used in this document is the EIF for the policy value in the
nonparametric model with censoring, up to the martingale augmentation noted after
Theorem~\ref{thm:ipcw}. That is a stronger reason to use it than double robustness alone.

\subsection{Neyman orthogonality}\label{app:neyman}

There is a compact way to say why the product-bias form of Theorem~\ref{thm:dr} is not a
coincidence.

\begin{definition}[Neyman orthogonality]\label{def:neyman}
Let $\psi(O;\eta)$ be an estimating score depending on nuisances $\eta$, with $\eta^\ast$ the
truth. The score is \emph{Neyman-orthogonal} if the Gateaux derivative of $\E[\psi(O;\eta)]$
with respect to $\eta$, evaluated at $\eta^\ast$, vanishes in every direction:
\[
\left.\frac{\partial}{\partial t}\,\E\bigl[\psi\bigl(O;\eta^\ast+t(\eta-\eta^\ast)\bigr)\bigr]\right|_{t=0}=0
\qquad\text{for all admissible }\eta .
\]
\end{definition}

\begin{proposition}[The augmented score is Neyman-orthogonal]\label{prop:neyman}
The score \eqref{eq:aipw} satisfies Definition~\ref{def:neyman} in $\eta=(m,e)$.
\end{proposition}

\begin{proof}
Take the path $m_t=m^\ast+t\,\Delta m$ and $e_t=e^\ast+t\,\Delta e$. Substituting into the
product-bias identity \eqref{eq:productbias} of Theorem~\ref{thm:dr},
\[
\E\bigl[\psi(O;m_t,e_t,\delta)\bigr]-\psi(P)
=\E\left[t\,\Delta m_{\delta}\left(1-\frac{e^\ast_\delta}{e^\ast_\delta+t\,\Delta e_\delta}\right)\right]
=\E\left[t\,\Delta m_{\delta}\cdot\frac{t\,\Delta e_\delta}{e^\ast_\delta+t\,\Delta e_\delta}\right],
\]
which is $t^2$ times a term bounded near $t=0$ by positivity. Differentiating at $t=0$ gives
zero.
\end{proof}

So the bias is second order in the nuisance error: first-order insensitivity is exactly what
permits slowly-converging nuisance estimators, and Corollary~\ref{cor:rates} is the quantitative
version. Orthogonality, double robustness, and the product-bias identity are three descriptions
of one property.

\subsection{Donsker classes, and why cross-fitting replaces them}\label{app:donsker}

One more piece of vocabulary appears in \S\ref{sec:crossfit}.

\begin{definition}[Donsker class]\label{def:donsker}
A class $\mathcal F$ of functions is $P$-Donsker if the empirical process
$f\mapsto n^{-1/2}\sum_i\{f(O_i)-\E f\}$ converges weakly, as a process indexed by
$\mathcal F$, to a tight Gaussian limit. Sufficient conditions restrict the complexity of
$\mathcal F$, for instance a finite bracketing-entropy integral.
\end{definition}

The classical route to asymptotic normality with estimated nuisances requires
$\hat\eta$ to lie, with probability tending to one, in a fixed Donsker class. That is a genuine
restriction: it holds for parametric and smoothly-constrained nonparametric estimators, and
fails for adaptively-tuned ensembles, whose effective complexity grows with the sample.

Proposition~\ref{prop:crossfit} shows the requirement is avoidable. By fitting $\hat\eta$ on
data independent of the fold where the score is evaluated, one conditions on $\hat\eta$ and the
empirical process degenerates to a sum of i.i.d.\ terms with a fixed function --- for which no
complexity control is needed at all. The cost is a constant factor in computation; the benefit
is that any consistent nuisance estimator becomes admissible.

%=====================================================================
\section{Causal identification}\label{app:causal}
%=====================================================================

Let $O=(X,A,Y)$ be one observation and let Assumptions~\ref{ass:consistency}--\ref{ass:positivity}
hold. Write $e_a(x)=\Prob(A=a\mid X=x)$ and $m_a(x)=\E[Y\mid X=x,A=a]$.

\begin{theorem}[Adjustment formula]\label{thm:adjust}
For any fixed $a$, $\ \E[Y(a)]=\E\bigl[\E[Y\mid X,A=a]\bigr]$, so $m_a(x)$ identifies
$\E[Y(a)\mid X=x]$.
\end{theorem}

\begin{proof}
Fix $x$ with $a\in\mathcal A(x)$; positivity guarantees $\Prob(A=a\mid X=x)>0$, so the
conditional expectation is well defined. Then
\[
\E[Y\mid X=x,A=a]
\;\overset{(i)}{=}\;\E[Y(a)\mid X=x,A=a]
\;\overset{(ii)}{=}\;\E[Y(a)\mid X=x],
\]
where $(i)$ is consistency --- on the event $A=a$ the observed $Y$ \emph{is} $Y(a)$, so the two
random variables agree pointwise there --- and $(ii)$ is unconfoundedness, since given $X$ the
event $A=a$ is independent of $Y(a)$ and property 3 of Appendix~\ref{app:condexp} applies.
Taking expectations over $X$ and using the tower property,
\[
\E[Y(a)]=\E\bigl[\E[Y(a)\mid X]\bigr]=\E\bigl[\E[Y\mid X,A=a]\bigr].\qedhere
\]
\end{proof}

\begin{corollary}[Value of a deterministic rule]\label{cor:value}
For any $\delta\in\mathcal D$, $\ \E[Y(\delta(X))]=\E\bigl[m_{\delta(X)}(X)\bigr]$.
\end{corollary}

\begin{proof}
Condition on $X=x$. Since $\delta$ is a deterministic function, $\delta(x)$ is a fixed element of
$\mathcal A(x)$, so $\E[Y(\delta(X))\mid X=x]=\E[Y(\delta(x))\mid X=x]=m_{\delta(x)}(x)$ by the
display in the proof of Theorem~\ref{thm:adjust}. The tower property finishes it.
\end{proof}

\begin{theorem}[Inverse-probability representation]\label{thm:ipw}
For any $\delta\in\mathcal D$,
$\ \E[Y(\delta(X))]=\E\left[\dfrac{\one{A=\delta(X)}}{e_{\delta(X)}(X)}\,Y\right]$.
\end{theorem}

\begin{proof}
Let $W=\one{A=\delta(X)}/e_{\delta(X)}(X)$ and condition on $X$. Since $\delta(X)$ and
$e_{\delta(X)}(X)$ are $X$-measurable, property 2 of Appendix~\ref{app:condexp} lets us pull the
denominator out:
\begin{align*}
\E[WY\mid X]
&= \frac{1}{e_{\delta(X)}(X)}\sum_{a\in\mathcal A(X)}\Prob(A=a\mid X)\,\one{a=\delta(X)}\,\E[Y\mid X,A=a]\\
&= \frac{\Prob(A=\delta(X)\mid X)}{e_{\delta(X)}(X)}\ \E[Y\mid X,A=\delta(X)]
\;=\; m_{\delta(X)}(X),
\end{align*}
because only the term $a=\delta(X)$ survives the indicator and the ratio is $1$ by definition of
$e$. Positivity bounds the denominator away from zero, so $WY$ is integrable. Taking expectations
over $X$ and applying Corollary~\ref{cor:value} gives the claim.
\end{proof}

Theorems~\ref{thm:adjust} and \ref{thm:ipw} give two identification routes with disjoint failure
modes: the first needs $m$, the second needs $e$. The next result combines them so that only one
need be correct.

\begin{definition}[Augmented score]\label{def:aipw}
For candidate nuisance functions $m$ and $e$ --- not necessarily the true ones ---
\begin{equation}
\psi(O;m,e,\delta)
= m_{\delta(X)}(X)
+ \frac{\one{A=\delta(X)}}{e_{\delta(X)}(X)}\Bigl(Y-m_{\delta(X)}(X)\Bigr).
\label{eq:aipw}
\end{equation}
\end{definition}

\begin{theorem}[Double robustness]\label{thm:dr}
Let $m^\ast,e^\ast$ be the true nuisance functions. For any $m$ and any $e$ bounded below,
\begin{equation}
\E\bigl[\psi(O;m,e,\delta)\bigr]-\E[Y(\delta(X))]
=\E\left[\Bigl(m_{\delta(X)}(X)-m^\ast_{\delta(X)}(X)\Bigr)
\left(1-\frac{e^\ast_{\delta(X)}(X)}{e_{\delta(X)}(X)}\right)\right].
\label{eq:productbias}
\end{equation}
Consequently the score is unbiased for the policy value if \emph{either} $m=m^\ast$ \emph{or}
$e=e^\ast$.
\end{theorem}

\begin{proof}
Abbreviate $m_\delta=m_{\delta(X)}(X)$ and similarly for $m^\ast_\delta,e_\delta,e^\ast_\delta$;
all four are $X$-measurable. Condition on $X$ and evaluate the augmentation term, using property
2 of Appendix~\ref{app:condexp} to extract $1/e_\delta$ and $m_\delta$:
\begin{align*}
\E\left[\frac{\one{A=\delta}}{e_\delta}\bigl(Y-m_\delta\bigr)\Bigm| X\right]
&= \frac{1}{e_\delta}\ \E\bigl[\one{A=\delta}\,(Y-m_\delta)\bigm| X\bigr]\\
&= \frac{1}{e_\delta}\ \Prob(A=\delta\mid X)\ \E\bigl[Y-m_\delta\bigm| X, A=\delta\bigr]\\
&= \frac{e^\ast_\delta}{e_\delta}\ \bigl(m^\ast_\delta-m_\delta\bigr),
\end{align*}
where the last step uses $\E[Y\mid X,A=\delta]=m^\ast_\delta$ (Theorem~\ref{thm:adjust}). Hence
\[
\E[\psi\mid X]
= m_\delta+\frac{e^\ast_\delta}{e_\delta}\bigl(m^\ast_\delta-m_\delta\bigr)
= m^\ast_\delta+\bigl(m_\delta-m^\ast_\delta\bigr)\left(1-\frac{e^\ast_\delta}{e_\delta}\right),
\]
the second equality being algebra. Taking expectations over $X$ and subtracting
$\E[Y(\delta(X))]=\E[m^\ast_\delta]$ from Corollary~\ref{cor:value} gives \eqref{eq:productbias}.
If $m=m^\ast$ the first factor of the integrand vanishes identically; if $e=e^\ast$ the second
does.
\end{proof}

\begin{corollary}[Rate double robustness]\label{cor:rates}
Suppose $\hat e\ge\epsilon$ and let $r_m,r_e$ be sequences such that, with probability tending to
one, $\|\hat m-m^\ast\|_{L_2(P)}\le r_m$ and $\|\hat e-e^\ast\|_{L_2(P)}\le r_e$. Then on that
event the bias \eqref{eq:productbias} is at most $\epsilon^{-1}r_mr_e$ in absolute value.
The bias is therefore negligible after the $\sqrt n$ scaling of \eqref{eq:asymplinear} whenever
$\sqrt n\,r_mr_e\to0$, that is, whenever the product of the two rates vanishes faster than
$n^{-1/2}$. It suffices, in particular, that each nuisance converge slightly faster than
$n^{-1/4}$; the exact rate $n^{-1/4}$ for both leaves the scaled bias bounded but not vanishing,
so a little margin is needed.
\end{corollary}

\begin{proof}
Under the lower bound $e\ge\epsilon$ we have $|1-e^\ast/e|=|e-e^\ast|/e\le|e-e^\ast|/\epsilon$.
Cauchy--Schwarz (Appendix~\ref{app:prereq}, item 2) applied to \eqref{eq:productbias} gives
$|\E[\psi]-\E[Y(\delta(X))]|\le\epsilon^{-1}\|\hat m-m^\ast\|_{L_2}\|\hat e-e^\ast\|_{L_2}$.
\end{proof}

\begin{remark*}
The $n^{-1/4}$ threshold is the whole practical content of the theorem, and it is worth stating
without symbols: each nuisance may be estimated at a rate far slower than the $n^{-1/2}$ demanded
of the target, because only the \emph{product} of the two errors enters the bias. That is the
licence for the flexible estimators of \S\ref{sec:models}, and it is why a modest improvement in
one nuisance can be traded against a deterioration in the other.
\end{remark*}

\begin{proposition}[Cross-fitting]\label{prop:crossfit}
Partition the sample into $K$ folds $I_1,\dots,I_K$, fit $\hat\eta^{(-k)}$ on the data outside
$I_k$, and form \eqref{eq:crossfit}. Then the empirical-process remainder
$|I_k|^{-1/2}\sum_{i\in I_k}\bigl\{f(O_i)-\E f\bigr\}$, with
$f=\psi(\cdot;\hat\eta^{(-k)})-\psi(\cdot;\eta^\ast)$, converges to zero in probability under
consistency of $\hat\eta^{(-k)}$ alone --- no complexity restriction on the estimator class is
required.
\end{proposition}

\begin{proof}
Condition on the out-of-fold data. Then $f$ is a fixed deterministic function, and the
$\{O_i\}_{i\in I_k}$ are i.i.d.\ draws from $P$ independent of it. Hence the conditional mean of
the display is zero, and its conditional variance is
$\Var\bigl(f(O)\bigr)\le\E[f(O)^2]=\|f\|_{L_2(P)}^2$, which converges to zero in probability as
$\hat\eta^{(-k)}\to\eta^\ast$. Chebyshev's inequality (Appendix~\ref{app:prereq}, item 4) gives,
for any $t>0$, $\Prob\bigl(|\,\cdot\,|>t\mid\text{out-of-fold}\bigr)\le\|f\|^2_{L_2(P)}/t^2$,
which tends to zero in probability; bounded convergence removes the conditioning. Averaging over
the $K$ folds preserves the conclusion.
\end{proof}

\begin{remark*}
Without cross-fitting, $f$ depends on the same observations at which it is evaluated, the
conditional-mean-zero step fails, and one must instead control $\sup_{\eta\in\mathcal H}|\cdot|$
over a class $\mathcal H$ containing $\hat\eta$ --- which is what a Donsker condition buys, and
what ensembles do not satisfy.
\end{remark*}

%=====================================================================
\section{Censoring}\label{app:censor}
%=====================================================================

\begin{lemma}[Restricted mean as an area]\label{lem:rmst-area}
For $T\ge0$ and $L>0$, $\ \E[\min(T,L)]=\int_0^L\Prob(T>t)\,dt$.
\end{lemma}

\begin{proof}
For fixed $\omega$, $\int_0^L\one{T(\omega)>t}\,dt$ equals $T(\omega)$ if $T(\omega)\le L$ and
$L$ otherwise, i.e.\ exactly $\min(T(\omega),L)$. The integrand is nonnegative and jointly
measurable, so Tonelli (Appendix~\ref{app:prereq}, item 1) permits exchanging the integral with
the expectation:
$\E[\min(T,L)]=\E\int_0^L\one{T>t}\,dt=\int_0^L\Prob(T>t)\,dt$.
\end{proof}

Let $D$ be the censoring time, $Y=\min(T,L)$, $R=\one{D\ge Y}$ the indicator that $Y$ is
observed, and $G(t\mid x,a)=\Prob(D\ge t\mid X=x,A=a)$.

\begin{theorem}[Inverse-probability-of-censoring weighting]\label{thm:ipcw}
Under Assumptions~\ref{ass:consistency}--\ref{ass:censor}, for any $a$,
\[
\E[Y(a)]=\E\left[\frac{R\ \one{A=a}}{e_a(X)\ G\bigl(Y\mid X,a\bigr)}\ Y\right].
\]
\end{theorem}

\begin{proof}
First recover the complete-data conditional mean. Condition on $(X,A,Y)$. Because $Y$ is a
function of $T$ and $D\indep T\mid(X,A)$ by Assumption~\ref{ass:censor},
\[
\E[R\mid X,A,Y]=\Prob(D\ge Y\mid X,A,Y)=\Prob(D\ge t\mid X,A)\Big|_{t=Y}=G(Y\mid X,A),
\]
where substituting the realised value of $Y$ is legitimate precisely because $D$ is
conditionally independent of $T$ and hence of $Y$. Since $Y\le L$, Assumption~\ref{ass:censor}
gives $G(Y\mid X,A)\ge\epsilon_D>0$, so the ratio below is integrable. Using property 2 of
Appendix~\ref{app:condexp} to extract the $(X,A,Y)$-measurable factor $Y/G(Y\mid X,A)$,
\[
\E\left[\frac{R\,Y}{G(Y\mid X,A)}\Bigm| X,A\right]
=\E\left[\frac{Y}{G(Y\mid X,A)}\,\E[R\mid X,A,Y]\Bigm| X,A\right]
=\E[Y\mid X,A].
\]
The left-hand side is a functional of observed data alone. Substituting it for
$\E[Y\mid X,A=a]$ in Theorem~\ref{thm:ipw} yields the statement.
\end{proof}

\begin{remark*}[Efficiency]
The score \eqref{eq:score} is the efficient influence function for the policy value in the
censored model \emph{if} one additionally augments over the censoring martingale: an extra term
that accumulates $G(u)^{-1}\E[Y\mid T>u,X,A]$ against the increments of the censoring
counting-process martingale $M^D$ over $u\in[0,L]$, correcting for misspecification of $G$ at
each point along the way rather than only at the observed exit time. Omitting it costs
efficiency, not consistency, when $G$ is correct. Since censoring here
is administrative and $G$ nearly known, the omission is defensible.
\end{remark*}

%=====================================================================
\section{Estimation lemmas}\label{app:estlemmas}
%=====================================================================

\subsection{Monotonicity of the constrained rule}\label{app:mono}

Recall $F_\tau(x)=\{a\in\mathcal A(x):p_a(x)\ge\tau\}$ and $\delta_\tau$ from
\eqref{eq:pointwise}.

\begin{proof}[Proof of Proposition~\ref{prop:mono}]
Fix $x$ and take $\tau'>\tau$, so $F_{\tau'}(x)\subseteq F_\tau(x)$.

\emph{Correctness rises.} Suppose $F_{\tau'}(x)\ne\emptyset$ and write $a_\tau=\delta_\tau(x)$,
$a_{\tau'}=\delta_{\tau'}(x)$. Either $a_\tau\in F_{\tau'}(x)$, in which case $a_\tau$ minimises
$m$ over $F_\tau(x)$ and lies in the subset $F_{\tau'}(x)$, hence also minimises $m$ over
$F_{\tau'}(x)$, so $a_{\tau'}=a_\tau$ up to ties and the two $p$ values coincide; or
$a_\tau\notin F_{\tau'}(x)$, in which case $p_{a_\tau}(x)<\tau'\le p_{a_{\tau'}}(x)$. Either way
$p_{\delta_{\tau'}(x)}(x)\ge p_{\delta_\tau(x)}(x)$. If $F_{\tau'}(x)=\emptyset$ the convention
selects the maximiser of $p$ over $\mathcal A(x)$, which dominates every alternative.

\emph{Duration rises.} Minimising a fixed function over a smaller set gives a weakly larger
minimum:
$m_{\delta_{\tau'}(x)}(x)=\min_{a\in F_{\tau'}(x)}m_a(x)\ge\min_{a\in F_\tau(x)}m_a(x)=m_{\delta_\tau(x)}(x)$,
and when $F_{\tau'}(x)=\emptyset$ the $p$-maximiser is weakly slower still.

Both inequalities hold pointwise in $x$; monotonicity of the expectation, together with
Corollary~\ref{cor:value} applied to $Y$ and to $C$, transfers them to $V_T$ and $V_C$.
\end{proof}

\begin{proof}[Proof of Corollary~\ref{cor:feasible}]
Feasibility is the defining property of the set in \eqref{eq:taustar}. By
Proposition~\ref{prop:mono}, $V_C(\delta_\tau)$ is nondecreasing, so the feasible grid points
form an upper set $\{\tau\in\mathcal T:\tau\ge\tau^\star\}$. Also by
Proposition~\ref{prop:mono}, $V_T(\delta_\tau)$ is nondecreasing, so among those points $V_T$ is
smallest at the smallest, namely $\tau^\star$.
\end{proof}

\subsection{Retransformation}

\begin{lemma}[Retransformation bias]\label{lem:jensen-retrans}
Let $\log(1+T)=\mu(x)+\varepsilon$ with $\E[\varepsilon\mid x]=0$. Then
$\E[1+T\mid x]=e^{\mu(x)}\E[e^\varepsilon\mid x]\ge e^{\mu(x)}$, with strict inequality unless
$\varepsilon$ is degenerate.
\end{lemma}

\begin{proof}
$1+T=e^{\mu(x)+\varepsilon}$, and $e^{\mu(x)}$ is $x$-measurable, so property 2 of
Appendix~\ref{app:condexp} gives $\E[1+T\mid x]=e^{\mu(x)}\E[e^\varepsilon\mid x]$. Jensen's
inequality (Appendix~\ref{app:prereq}, item 3) applied to the convex map $u\mapsto e^u$ gives
$\E[e^\varepsilon\mid x]\ge e^{\E[\varepsilon\mid x]}=e^0=1$, with equality only if
$\varepsilon$ is a.s.\ constant, since $u\mapsto e^u$ is strictly convex.
\end{proof}

\subsection{The winner's curse}

\begin{lemma}[Optimistic bias of a plug-in minimum]\label{lem:winner}
For any integrable random variables $\hat m_1,\dots,\hat m_K$,
$\ \E\bigl[\min_a\hat m_a\bigr]\le\min_a\E[\hat m_a]$.
\end{lemma}

\begin{proof}
For each fixed index $a_0$, $\min_a\hat m_a\le\hat m_{a_0}$ pointwise, so monotonicity of the
expectation gives $\E[\min_a\hat m_a]\le\E[\hat m_{a_0}]$. As this holds for every $a_0$ it
holds for the minimising one. (Equivalently: $(u_1,\dots,u_K)\mapsto\min_au_a$ is concave, and
Jensen's inequality for concave functions applies.)
\end{proof}

\subsection{Effective sample size}

\begin{lemma}[Kish]\label{lem:ess}
Let $Y_1,\dots,Y_n$ be i.i.d.\ with variance $\sigma^2$ and let $w_i\ge0$ be fixed weights, not
all zero. Then $\Var\bigl(\sum_iw_iY_i/\sum_iw_i\bigr)=\sigma^2/\text{ESS}$ with
$\text{ESS}=(\sum_iw_i)^2/\sum_iw_i^2$.
\end{lemma}

\begin{proof}
By independence, $\Var(\sum_iw_iY_i)=\sigma^2\sum_iw_i^2$; dividing by the square of the
constant $\sum_iw_i$ gives $\sigma^2\sum_iw_i^2/(\sum_iw_i)^2$. Setting this equal to
$\sigma^2/n_{\text{eff}}$ and solving identifies $n_{\text{eff}}=\text{ESS}$. Cauchy--Schwarz
gives $\text{ESS}\le n$, with equality iff all weights are equal.
\end{proof}

\subsection{Clustered variance}

\begin{lemma}[Design effect]\label{lem:cluster}
Suppose the sample consists of $G$ clusters of common size $\bar n_g$, that
$\Var(\psi_i)=\sigma^2$ for all $i$, and that $\Corr(\psi_i,\psi_j)=\rho$ for $i\ne j$ in the
same cluster and $0$ across clusters. Then
$\Var(\bar\psi)=\frac{\sigma^2}{n}\bigl[1+(\bar n_g-1)\rho\bigr]$.
\end{lemma}

\begin{proof}
Write $n=G\bar n_g$. Then
$\Var\bigl(\sum_i\psi_i\bigr)=\sum_i\Var(\psi_i)+\sum_{i\ne j}\Cov(\psi_i,\psi_j)
=n\sigma^2+G\,\bar n_g(\bar n_g-1)\rho\sigma^2=n\sigma^2[1+(\bar n_g-1)\rho]$,
since each cluster contributes $\bar n_g(\bar n_g-1)$ ordered within-cluster pairs and
cross-cluster covariances vanish. Dividing by $n^2$ gives the claim.
\end{proof}

The bracketed factor is the design effect. With $\bar n_g$ in the hundreds, even $\rho=0.01$
roughly doubles the variance, which is why \S\ref{sec:variance} clusters by district--year.

%=====================================================================
\section{The marginal sensitivity model}\label{app:sens}
%=====================================================================

\begin{definition}[Marginal sensitivity model]\label{def:msm}
Let $e^\ast_a(x)$ be the propensity given $X$ and $\tilde e_a(x,u)$ the propensity given $X$ and
an unobserved $U$. For $\Gamma\ge1$ the model asserts
\[
\frac1\Gamma\ \le\
\frac{\tilde e_a(x,u)\big/\bigl(1-\tilde e_a(x,u)\bigr)}
     {e^\ast_a(x)\big/\bigl(1-e^\ast_a(x)\bigr)}
\ \le\ \Gamma
\qquad\text{for all }x,u,a.
\]
\end{definition}

$\Gamma$ reads directly: an unmeasured factor may change the odds of receiving a given flow by at
most a factor $\Gamma$, beyond everything in $X$. $\Gamma=1$ is Assumption~\ref{ass:unconf}.

\begin{proposition}[Extremal bounds]\label{prop:msm-bounds}
Write the true weight as $\tilde w=1/\tilde e_{\delta(X)}(X,U)$. Under Definition~\ref{def:msm}
each $\tilde w_i$ is confined to an interval $[\underline w_i,\overline w_i]$ obtained by
inverting the odds-ratio bounds, and the sharp bounds on the policy value solve
\[
\overline V(\Gamma)=\max_{\tilde w\in\mathcal W_\Gamma}
\frac{\sum_i\tilde w_i\one{A_i=\delta(X_i)}Y_i}{\sum_i\tilde w_i\one{A_i=\delta(X_i)}},
\qquad
\underline V(\Gamma)=\min_{\tilde w\in\mathcal W_\Gamma}(\cdot),
\]
with $\mathcal W_\Gamma=\prod_i[\underline w_i,\overline w_i]$. Each is a linear-fractional
program, attains its optimum at a vertex of the box, and the optimal vertex assigns
$\overline w_i$ to observations with $Y_i$ above the optimised value and $\underline w_i$ to the
rest.
\end{proposition}

\begin{proof}[Proof sketch]
The pointwise interval follows by inverting the odds-ratio bounds of Definition~\ref{def:msm}.
For the program, fix a candidate value $v$ and note that
$\sum_i\tilde w_i\one{\cdot}(Y_i-v)\ge0$ iff the weighted mean is at least $v$; the left side is
linear in $\tilde w$, so its maximum over a box is attained by setting $\tilde w_i=\overline w_i$
when $Y_i>v$ and $\underline w_i$ otherwise. Monotonicity of this maximum in $v$ makes the
optimal $v$ the unique fixed point, found by sorting on $Y_i$ once and scanning the $n+1$
candidate thresholds in a single pass, so the cost is that of the sort.
\end{proof}

The reportable summary is the \emph{breakdown point}
$\Gamma^\dagger=\inf\{\Gamma:0\in[\underline\Delta(\Gamma),\overline\Delta(\Gamma)]\}$, the
strength of unmeasured confounding at which the conclusion first fails.

\end{document}
